\documentclass[11pt,letterpaper]{article}

% ── Packages ──
\usepackage[margin=1in,headheight=14pt]{geometry}
\usepackage{amsmath,amssymb,amsthm}
\usepackage{mathtools}
\usepackage{enumitem}
\usepackage{hyperref}
\usepackage{cleveref}
\usepackage{tocloft}
\usepackage{fancyhdr}
\usepackage{titlesec}
\usepackage{tikz-cd}
\usepackage{booktabs}

% ── Theorem environments ──
\newtheorem{conjecture}{Conjecture}[section]
\newtheorem{theorem}[conjecture]{Theorem}
\newtheorem{corollary}[conjecture]{Corollary}
\newtheorem{lemma}[conjecture]{Lemma}
\newtheorem{definition}[conjecture]{Definition}
\newtheorem{proposition}[conjecture]{Proposition}
\newtheorem{example}[conjecture]{Example}
\newtheorem{remark}[conjecture]{Remark}

% ── Header ──
\pagestyle{fancy}
\fancyhf{}
\lhead{B.\ Davis}
\rhead{Renormalization of Completion}
\cfoot{\thepage}

% ── Commands (inherited from Branches VI--VII) ──
\newcommand{\Hol}{\mathrm{Hol}}
\newcommand{\Khat}{\hat{K}}
\newcommand{\tbudget}{\tau_{\mathrm{budget}}}
\newcommand{\smax}{s_{\mathrm{max}}}
\newcommand{\Svalid}{S_{\mathrm{valid}}}
\newcommand{\tr}{\operatorname{tr}}
\newcommand{\Vol}{\operatorname{Vol}}
\newcommand{\Cech}{\check{C}}
\newcommand{\Hcech}{\check{H}}
\newcommand{\sheaf}{\mathcal{F}}
\newcommand{\cover}{\mathcal{U}}
\newcommand{\nerve}{\mathcal{N}}
\newcommand{\constraint}{\mathcal{C}}
\DeclareMathOperator{\im}{im}
\DeclareMathOperator{\coker}{coker}
\DeclareMathOperator{\Hom}{Hom}
\DeclareMathOperator{\Ext}{Ext}
\DeclareMathOperator{\rk}{rk}
\DeclareMathOperator{\diam}{diam}

% ── New commands for Branch VIII ──
\newcommand{\block}{\mathcal{B}}
\newcommand{\partition}{\mathcal{P}}
\newcommand{\SO}{\mathrm{SO}}

\begin{document}

\title{\textbf{Renormalization of Completion}\\[6pt]
\large Renormalization Group Flow on the Davis Manifold\\[4pt]
\normalsize Branch VIII: Coarse-Graining, Fixed Points, and the\\
Scale-Dependent Structure of Completion}
\author{Bee Rosa Davis}
\date{February 2026}
\maketitle

\begin{abstract}
We construct a renormalization group (RG) flow on the space of 
completion problems by defining a coarse-graining operation on the 
nerve of the constraint cover. Merging nearby constraint patches into 
blocks and pushing forward the completion presheaf defines a 
semigroup of scale transformations $\{T_\ell\}_{\ell \geq 1}$ acting 
on pairs $(\nerve, \sheaf)$. The trichotomy parameter $\Gamma$ of the 
parent framework becomes scale-dependent: $\Gamma(\ell)$ flows under 
coarse-graining, and its beta function $\beta(\Gamma) = 
d\Gamma / d(\ln \ell)$ has exactly three fixed points corresponding 
to the determined ($\Gamma^* = \infty$), critical ($\Gamma^* = 1$), 
and underdetermined ($\Gamma^* = 0$) regimes. The critical fixed 
point is unstable, with the relevant eigenvalue $y_t = 1/\nu$ 
recovering the correlation length exponent of Branch~VI. A 
C-theorem analog (monotonically decreasing completion entropy) 
follows from a Cauchy interlace inequality on the coarsened \v{C}ech 
Laplacian. Cross-synthesis with the parent framework interprets 
transformer layers as successive RG steps, predicting scale-dependent 
completion capacity and connecting semantic depth to the number of 
coarse-graining steps needed to reach the determined fixed point.

\medskip\noindent
\textbf{Result count:} 8 foundational results (RG1--RG8), 7 
first-order derivations, 11 cross-synthesis points with the parent 
framework and Branches~VI--VII. Total: 26 new results (4 theorems, 
22 conjectures).
\end{abstract}

\tableofcontents
\newpage

% ═══════════════════════════════════════════════════════
\section{Motivation: Why Renormalization?}
% ═══════════════════════════════════════════════════════

Branch~VI establishes that the completion transition at $\Gamma = 1$ 
has the hallmarks of a genuine phase transition: divergent specific 
heat, divergent correlation length, and critical exponents satisfying 
scaling relations. Branch~VII identifies the \v{C}ech Laplacian as 
the natural operator governing these phenomena, with the spectral 
gap $\lambda_1(\Delta_1)$ playing the role of the inverse squared 
correlation length.

What is missing is a \emph{scale-dependent} theory. In the current 
framework, the trichotomy parameter $\Gamma$ and the curvature 
$\Khat$ are properties of the problem as a whole. But real 
completion problems --- and especially transformer-based language 
models --- process information \emph{hierarchically}: local 
constraints are resolved first, and global coherence emerges through 
successive layers of integration. This multi-scale structure demands 
a renormalization group.

The RG program in condensed matter physics asks three questions:

\begin{enumerate}
    \item \textbf{Coarse-graining:} How do you systematically 
    eliminate fine-scale degrees of freedom while preserving 
    long-range structure?
    
    \item \textbf{Fixed points:} What are the scale-invariant 
    theories that anchor the space of all theories?
    
    \item \textbf{Universality:} Why do microscopically different 
    systems share the same critical behavior?
\end{enumerate}

The Davis framework contains natural answers to all three:

\begin{enumerate}
    \item Coarse-graining means merging constraint patches on the 
    nerve and pushing forward the completion presheaf (a 
    block-spin transformation on the nerve).
    
    \item The three fixed points are the three regimes: determined, 
    critical, and underdetermined.
    
    \item Universality arises because the critical exponents depend 
    only on the Hausdorff dimension of the constraint graph and 
    the symmetry group of the constraint structure, not on the 
    microscopic details of the presheaf.
\end{enumerate}

This branch makes all three connections precise.

% ═══════════════════════════════════════════════════════
\section{The RG Dictionary}
% ═══════════════════════════════════════════════════════

\begin{center}
\renewcommand{\arraystretch}{1.4}
\begin{tabular}{lll}
\textbf{Statistical Mechanics} & \textbf{Davis Framework} 
& \textbf{Equation} \\
\hline
Lattice sites & Nerve vertices (constraints) 
& $v \in \nerve^{(0)}$ \\
Block-spin partition $\partition_\ell$ & Constraint grouping 
& $\block_a \subset \nerve^{(0)}$ \\
Block-spin variable & Coarsened presheaf section 
& $s_a \in \sheaf_\ell(\block_a)$ \\
Coarsened lattice & Coarsened nerve $\nerve_\ell$ 
& $|\nerve_\ell^{(0)}| < |\nerve^{(0)}|$ \\
Effective Hamiltonian $H_{\mathrm{eff}}$ & Coarsened Laplacian 
$\Delta_0^{(\ell)}$ & Leading-order RG \\
Mass gap $m$ & Spectral gap $\lambda_1(\Delta_1^{(\ell)})$ 
& $m^2 = \lambda_1$ \\
Running coupling $g(\ell)$ & Scale-dependent $\Gamma(\ell)$ 
& $\beta(\Gamma)$ flow \\
UV fixed point & Fine-scale completion problem 
& $\ell = 0$ \\
IR fixed point & Coarsest completion problem 
& $\ell = \ell_{\max}$ \\
Relevant operator & Obstruction surviving coarse-graining 
& $[\alpha] \in \Hcech^1(\nerve_\ell, \sheaf_\ell)$ \\
Irrelevant operator & Sub-scale obstruction 
& $[\alpha] \in \ker(\Hcech^1 \to \Hcech^1_\ell)$ \\
C-theorem (Zamolodchikov) & Completion entropy monotonicity 
& $C(\ell_2) \leq C(\ell_1)$ \\
\end{tabular}
\end{center}

\begin{remark}[What is Being Coarse-Grained]
A critical distinction: the coarse-graining operates on the 
\emph{nerve} $\nerve(\cover)$ of the constraint cover, whose 
vertices are \emph{constraints} (rows, columns, boxes in Sudoku), 
not on the manifold $M$ directly. The nerve is the combinatorial 
skeleton encoding which constraints interact. A block-spin step 
merges nearby \emph{constraints} into super-constraints, not 
nearby cells into super-cells. This is the natural scale 
transformation because the completion obstruction lives in the 
\v{C}ech cohomology of the nerve, not in the homology of $M$.
\end{remark}

\begin{remark}[Linearization Regime]\label{rem:rg_linearization}
All quantitative RG results in this branch (spectral gaps, 
partition functions, trace inequalities, anomalous dimensions) 
operate in the linearized regime established in Branch~VII: the 
presheaf $\sheaf^{\mathrm{lin}}$ with additive coboundary operators 
on $\mathfrak{gl}(W)$-valued cochains. The exact non-abelian 
$\mathrm{GL}(W)$-valued cocycle condition is preserved qualitatively 
(gauge triviality, fixed-point structure) but the spectral analysis 
requires the abelian setting. As established in Branch~VII, 
Remark~5.1 (BCH Analytic Control), the linearization is valid when 
$\|\alpha\|_{\Cech^1} < (\log 2) / L$ where $L$ is the nerve 
diameter. Coarse-graining \emph{reduces} the nerve diameter, so 
the linearized regime becomes \emph{more} accurate at larger scales 
(see Y6 below).
\end{remark}

\begin{example}[The Shidoku Nerve]\label{ex:shidoku_nerve}
A $4 \times 4$ Shidoku has 12 constraints: 4 rows ($r_1, \ldots, 
r_4$), 4 columns ($c_1, \ldots, c_4$), and 4 boxes ($b_1, \ldots, 
b_4$, each $2 \times 2$). The nerve $\nerve(\cover)$ has 12 vertices 
and 32 edges (16 row-column, 8 row-box, 8 column-box), with the 
adjacency structure determined by cell sharing.

A natural coarse-graining groups constraints that share a band:
\begin{align*}
    \block_1 &= \{r_1, r_2, b_1, b_2\} 
    && \text{(top rows + top boxes)} \\
    \block_2 &= \{r_3, r_4, b_3, b_4\} 
    && \text{(bottom rows + bottom boxes)} \\
    \block_3 &= \{c_1, c_2\} && \text{(left columns)} \\
    \block_4 &= \{c_3, c_4\} && \text{(right columns)}
\end{align*}
The coarsened nerve $\nerve_1$ has 4 vertices with edges between 
every row-band block and every column block, forming the complete 
bipartite graph $K_{2,2}$.
\end{example}

% ═══════════════════════════════════════════════════════
\newpage
\part{Foundational Results (RG1--RG8)}
% ═══════════════════════════════════════════════════════

% ── RG1 ──
\section{RG1: Block-Spin Coarse-Graining on the Nerve}

\begin{definition}[Coarse-Graining Triple]\label{def:coarse_grain}
Let $\nerve = \nerve(\cover)$ be the nerve of the constraint cover 
$\cover = \{U_c\}_{c \in \constraint}$, with graph metric $d_\nerve$ 
on the 1-skeleton.

A \emph{scale-$\ell$ coarse-graining} of $(\nerve, \sheaf)$ is a 
triple $(\nerve_\ell, \pi_\ell, \sheaf_\ell)$ defined as follows:

\begin{enumerate}[label=(\roman*)]
    \item \textbf{Partition.} $\partition_\ell = 
    \{\block_a\}_{a \in A}$ is a partition of the vertex set 
    $\nerve^{(0)}$ such that $\diam(\block_a) \leq \ell$ in 
    $d_\nerve$ for all blocks $\block_a$.
    
    \item \textbf{Coarsened nerve.} $\nerve_\ell$ is the simplicial 
    complex with vertex set $A$ and a $k$-simplex 
    $[a_0, \ldots, a_k]$ whenever there exist 
    $v_i \in \block_{a_i}$ forming a $k$-simplex in $\nerve$.
    
    \item \textbf{Quotient map.} $\pi_\ell : \nerve \to \nerve_\ell$ 
    is the induced surjective simplicial map.
    
    \item \textbf{Coarsened presheaf.} The pushed-forward completion 
    presheaf is the inverse limit:
    \begin{equation}\label{eq:coarsened_presheaf}
        \sheaf_\ell(\block_a) = \varprojlim_{v \in \block_a} 
        \sheaf(v)
    \end{equation}
    where the limit is taken over all pairwise overlaps within the 
    block. That is: a section of $\sheaf_\ell$ on $\block_a$ is a 
    family of local completions on every constraint in $\block_a$ 
    that are mutually compatible on all pairwise overlaps within the 
    block.
\end{enumerate}
\end{definition}

\begin{theorem}[Injectivity of Coarse-Graining]\label{thm:rg_inject}
The coarse-graining construction yields a well-defined presheaf 
$\sheaf_\ell$ on $\nerve_\ell$. There is a canonical restriction map 
on global sections:
\begin{equation}\label{eq:rg_restriction}
    r_\ell : \Gamma(\nerve, \sheaf) \longrightarrow 
    \Gamma(\nerve_\ell, \sheaf_\ell)
\end{equation}
and $r_\ell$ is injective: every global completion of the original 
problem induces a global completion of the coarsened problem.
\end{theorem}

\begin{proof}
\textbf{Well-definedness.} Each stalk $\sheaf_\ell(\block_a)$ is 
an inverse limit of sets (the local section spaces connected by 
restriction maps), which exists in the category of sets. The 
restriction maps of $\sheaf_\ell$ (for inclusions between simplices 
of $\nerve_\ell$) are induced by the restriction maps of $\sheaf$ 
on the corresponding sub-collections of the original nerve.

\textbf{Canonical map.} A global section $s \in \Gamma(\nerve, 
\sheaf)$ is a compatible family $(s_v)_{v \in \nerve^{(0)}}$. For 
each block $\block_a$, the restriction $(s_v)_{v \in \block_a}$ is 
automatically compatible on all pairwise overlaps within $\block_a$ 
(since it is compatible on \emph{all} overlaps globally). Therefore 
$(s_v)_{v \in \block_a} \in \sheaf_\ell(\block_a)$.

\textbf{Injectivity.} If $r_\ell(s) = r_\ell(s')$, then for every 
block $\block_a$ and every $v \in \block_a$, $s_v = s'_v$. Since 
the blocks partition $\nerve^{(0)}$, $s_v = s'_v$ for all $v$, 
hence $s = s'$.
\end{proof}

\begin{remark}[Non-Surjectivity and False Completions]
\label{rem:false_completions}
The map $r_\ell$ is generally \emph{not} surjective: the coarsened 
problem $(\nerve_\ell, \sheaf_\ell)$ can have global sections that 
do not lift to global sections of $(\nerve, \sheaf)$. These are 
\emph{false completions} --- configurations that look valid at scale 
$\ell$ (compatible within each block and between blocks) but fail at 
finer scales (incompatible within a block when one examines the 
internal structure). The obstruction to lifting is classified by the 
Leray spectral sequence (RG1a below): the higher direct images 
$R^q \pi_{\ell *} \sheaf$ for $q \geq 1$ measure the sub-scale 
obstructions invisible at scale $\ell$.

\textbf{Linearization scope.} The coarse-graining construction 
(Definition~\ref{def:coarse_grain}) operates on the set-valued 
presheaf $\sheaf$ and requires no linearization. The spectral 
analysis (RG3 onward: Laplacians, partition functions, spectral 
gaps, anomalous dimensions) requires the linearized presheaf 
$\sheaf^{\mathrm{lin}}$. The Leray spectral sequence 
(Theorem~\ref{thm:leray_rg}) is stated for the abelian presheaf 
$\sheaf^{\mathrm{lin}}$. See Remark~\ref{rem:rg_linearization} 
for the linearization regime.
\end{remark}

\begin{example}[Shidoku Coarse-Graining]
Continuing Example~\ref{ex:shidoku_nerve}: the partition 
$\partition_1 = \{\block_1, \block_2, \block_3, \block_4\}$ groups 
the 12 nerve vertices into 4 blocks. The stalk dimensions of the 
coarsened presheaf are:
\begin{itemize}
    \item $\sheaf_1(\block_1)$: all completions of the top two rows 
    that are simultaneously valid for both rows \emph{and} both 
    top boxes. This is a \emph{stronger} constraint than requiring 
    validity for each constraint individually.
    \item $\sheaf_1(\block_3)$: all completions of the two left 
    columns, constrained to be mutually compatible (which is 
    automatic since distinct columns share no cells in a $4 \times 4$ 
    grid, so $\sheaf_1(\block_3) \cong \sheaf(c_1) \times 
    \sheaf(c_2)$).
\end{itemize}
The coarsened nerve $\nerve_1 \cong K_{2,2}$ has 4 vertices, 4 
edges, and no higher simplices. Global sections of $\sheaf_1$ on 
$K_{2,2}$ are band-level compatible completions.
\end{example}

\textbf{Derived from:} Branch~VII, SH1 (completion presheaf) + SH2 
(\v{C}ech complex) + standard inverse limit construction.

\textbf{Interpretation:} This is the Kadanoff block-spin 
transformation transplanted from the square lattice to the nerve of 
the constraint cover. Each ``super-spin'' (block section) is a 
coarsened description of the constraint structure at scale $\ell$.

% ── RG2 ──
\section{RG2: The Coarse-Graining Semigroup}

\begin{definition}[Filtration of Partitions and RG 
Semigroup]\label{def:rg_semigroup}
A \emph{hierarchical coarse-graining} of $(\nerve, \sheaf)$ is a 
filtration of partitions:
\begin{equation}\label{eq:filtration}
    \partition_0 = \{\{v\}\}_{v \in \nerve^{(0)}} 
    \preccurlyeq \partition_1 \preccurlyeq \cdots 
    \preccurlyeq \partition_L = \{\nerve^{(0)}\}
\end{equation}
where $\partition_i \preccurlyeq \partition_j$ means every block of 
$\partition_j$ is a union of blocks of $\partition_i$ (the coarser 
partition refines the finer one). The \emph{scale} $\ell$ ranges 
from $0$ (trivial partition, original nerve) to $L$ (single block, 
maximally coarsened).

The filtration induces a directed system of simplicial maps:
\begin{equation}\label{eq:directed_system}
    \nerve = \nerve_0 \xrightarrow{\pi_{0,1}} \nerve_1 
    \xrightarrow{\pi_{1,2}} \cdots 
    \xrightarrow{\pi_{L-1,L}} \nerve_L = \{*\}
\end{equation}
with pushed-forward presheaves $\sheaf_0 = \sheaf, \sheaf_1, 
\ldots, \sheaf_L$.

The \emph{RG semigroup} $T = \{T_\ell\}_{\ell = 0}^{L}$ is defined 
by $T_\ell(\nerve, \sheaf) = (\nerve_\ell, \sheaf_\ell)$, acting on 
the space of completion problems 
$\mathcal{P} = \{(\nerve, \sheaf) : \nerve \text{ finite simplicial 
complex}, \sheaf \text{ presheaf on } \nerve\}$.

\textbf{Effective parameters.} Each completion problem 
$(\nerve_\ell, \sheaf_\ell)$ induces effective parameters at 
scale $\ell$:
\begin{align}
    n(\ell) &= |\nerve_\ell^{(0)}| 
    && \text{(number of effective constraints)} 
    \label{eq:n_ell}\\
    m(\ell)^2 &= \lambda_1(\Delta_1^{(\ell)}) 
    && \text{(spectral gap of degree-1 Laplacian)} 
    \label{eq:m_ell}\\
    s(\ell) &= \frac{1}{n(\ell)} \sum_a \dim \sheaf_\ell(\block_a) 
    && \text{(average stalk dimension)} 
    \label{eq:s_ell}
\end{align}
The \emph{scale-dependent trichotomy parameter} is defined 
operationally:
\begin{equation}\label{eq:Gamma_ell}
    \Gamma(\ell) = \text{the parent's } \Gamma \text{ evaluated 
    on the coarsened problem } 
    (\nerve_\ell, \sheaf_\ell)
\end{equation}
That is, $\Gamma(\ell)$ inherits whatever functional form relates 
$\Gamma$ to the curvature, constraint count, and tolerance in the 
parent framework, applied to the effective quantities at scale 
$\ell$. At $\ell = 0$, $\Gamma(0)$ is the original trichotomy 
parameter.
\end{definition}

\begin{remark}[Wilson vs.\ Kadanoff Convention]\label{rem:wilson}
The tolerance budget $\tau$ (a property of the model, not the 
problem) is held fixed under coarse-graining. This is Wilson's 
convention: fix the ``temperature'' and coarse-grain the 
``Hamiltonian'' (the constraint structure). An alternative 
Kadanoff convention rescales $\tau(\ell) = \tau_0 \cdot 
\ell^{y_\tau}$ for some scaling dimension $y_\tau$, which can be 
absorbed into a redefinition of $\Gamma(\ell)$. We adopt the Wilson 
convention throughout.
\end{remark}

\textbf{Derived from:} RG1 (coarse-graining triple).

\textbf{Interpretation:} The filtration of partitions is the analog 
of choosing a specific blocking scheme in real-space 
renormalization. The nesting requirement ensures that each RG step 
groups previously-grouped constraints, never splitting a block. 
This is the minimal structure needed for the semigroup property 
$T_{\ell_2} = T_{\ell_2, \ell_1} \circ T_{\ell_1}$ to hold.

\begin{remark}[Semigroup Functoriality]\label{rem:semigroup}
The semigroup property holds because the pushforward of presheaves 
is functorial: if $\pi_{\ell_1, \ell_2} : \nerve_{\ell_1} \to 
\nerve_{\ell_2}$ is the simplicial map induced by coarsening 
$\partition_{\ell_1}$ to $\partition_{\ell_2}$, then 
$(\pi_{\ell_1, \ell_2})_* \sheaf_{\ell_1} = \sheaf_{\ell_2}$. 
This follows from the universal property of inverse limits: the 
inverse limit defining $\sheaf_{\ell_2}(\block_a^{(\ell_2)})$ over 
all vertices $v \in \block_a^{(\ell_2)}$ factors through the 
intermediate inverse limits $\sheaf_{\ell_1}(\block_b^{(\ell_1)})$ 
for blocks $\block_b^{(\ell_1)} \subseteq \block_a^{(\ell_2)}$ 
(which exist because the partitions are nested). The nesting 
$\partition_{\ell_1} \preccurlyeq \partition_{\ell_2}$ is essential: 
without it, the factorization fails and the semigroup property 
need not hold.
\end{remark}

% ── RG3 ──
\section{RG3: The \v{C}ech Laplacian at Scale $\ell$}

\begin{definition}[Coarsened \v{C}ech 
Complex]\label{def:coarsened_cech}
At scale $\ell$, the \emph{coarsened \v{C}ech complex} is:
\begin{equation}\label{eq:coarsened_complex}
    0 \longrightarrow \Cech^0(\nerve_\ell, \sheaf_\ell^{\mathrm{lin}}) 
    \xrightarrow{\;\delta_\ell^0\;}
    \Cech^1(\nerve_\ell, \sheaf_\ell^{\mathrm{lin}}) 
    \xrightarrow{\;\delta_\ell^1\;} \cdots
\end{equation}
with coboundary operators $\delta_\ell^p$ defined by the restriction 
maps of $\sheaf_\ell^{\mathrm{lin}}$ (the linearization of the 
coarsened presheaf, in the additive regime of Branch~VII).

The \emph{coarsened \v{C}ech Laplacian} on $p$-cochains is:
\begin{equation}\label{eq:coarsened_laplacian}
    \Delta_p^{(\ell)} = P_\ell^{(p)\dagger} \, \Delta_p \, P_\ell^{(p)}
\end{equation}
where $P_\ell^{(p)} : V_\ell^{(p)} \hookrightarrow 
\Cech^p(\nerve, \sheaf^{\mathrm{lin}})$ is the inclusion of the 
subspace of \emph{block-compatible $p$-cochains} and 
$P_\ell^{(p)\dagger}$ is its adjoint (orthogonal projection).

\textbf{Inherited inner product.} A block-compatible 0-cochain is 
constant on each block $\block_a$; a block-compatible 1-cochain 
assigns the same value to all original edges connecting a given 
pair of blocks (and zero to within-block edges). The inner product 
on $V_\ell^{(p)}$ is inherited from the ambient 
$\Cech^p(\nerve, \sheaf^{\mathrm{lin}})$ by restriction.

\textbf{Equivalently,} $\Delta_p^{(\ell)}$ is the Laplacian of the 
coarsened coboundary operators $\delta_\ell^p$ with \emph{weighted} 
inner product on $\Cech^p(\nerve_\ell, \sheaf_\ell^{\mathrm{lin}})$: 
each coarsened $p$-simplex $(a_0, \ldots, a_p)$ carries weight 
equal to the number of original $p$-simplices it represents. This 
edge-multiplicity weighting ensures the coarsened coboundary 
structure faithfully represents the fine-scale spectral content.
\end{definition}

\begin{remark}[Compression vs.\ Coarsened Nerve 
Laplacian]\label{rem:compression_vs_coarse}
The operator $\Delta_p^{(\ell)} = P_\ell^{(p)\dagger} \Delta_p 
P_\ell^{(p)}$ is \emph{not} the bare graph Laplacian of the 
coarsened nerve $\nerve_\ell$ with unit edge weights. That ``bare'' 
Laplacian forgets how many original edges each coarsened edge 
represents, and has a different spectrum. The C-theorem 
(Theorem~\ref{thm:c_theorem}) and Cauchy interlace both require the 
\emph{compression} definition, which preserves the inclusion 
$V_{\ell_2} \subseteq V_{\ell_1}$ and the associated spectral 
ordering. For the Shidoku nerve with the block partition of 
Appendix~\ref{app:shidoku_rg}, the compression has eigenvalues 
$\{0, 3, 6, 9\}$ while the bare coarsened Laplacian has 
$\{0, 2, 2, 4\}$ --- the C-theorem applies to the former.
\end{remark}

\begin{conjecture}[Spectral Gap Scaling]\label{conj:spectral_gap}
The spectral gap $m(\ell)^2 = \lambda_1(\Delta_1^{(\ell)})$ 
satisfies the scaling relation:
\begin{equation}\label{eq:spectral_scaling}
    m(\ell)^2 = \ell^{-2} \cdot m_R(g(\ell))^2
\end{equation}
where $m_R$ is the \emph{renormalized mass} and $g(\ell) = 
\Gamma(\ell)$ is the running coupling. At the critical point 
($\Gamma = 1$), $m_R = 0$ and the spectrum is gapless at all 
scales (scale invariance). Away from criticality, 
$m_R \neq 0$ and $m(\ell)$ vanishes at the scale $\ell = \xi$ 
(the correlation length).

The $\ell^{-2}$ prefactor is the standard engineering dimension 
of the mass gap under block-spin RG: since blocks of diameter 
$\ell$ merge $\sim\!\ell^{d_H}$ vertices, the effective spectral 
gap (an eigenvalue of a matrix whose dimension decreases with 
$\ell$) generically increases as $\ell^{-2}$ absent anomalous 
corrections. For discrete nerves, this should be understood as 
the leading scaling behavior of the spectral gap, to be verified 
numerically in specific examples (Appendix~A).
\end{conjecture}

\textbf{Derived from:} Branch~VII, SH7 (Euler characteristic, 
Laplacian) + X3 (correlation length from spectral gap).

\textbf{Interpretation:} The spectral gap of $\Delta_1^{(\ell)}$ 
is the inverse squared correlation length at scale $\ell$. When 
the gap vanishes, the system has long-range correlations at that 
scale --- the hallmark of criticality. The scaling relation 
connects the bare spectral gap to the renormalized mass through 
the running coupling.

% ── RG4 ──
\section{RG4: The Beta Function}

\begin{conjecture}[Beta Function of the Trichotomy 
Parameter]\label{conj:beta}
The RG flow of $\Gamma$ is governed by the beta function:
\begin{equation}\label{eq:beta_def}
    \beta(\Gamma) := \frac{d\Gamma}{d(\ln \ell)}
\end{equation}
For generic completion problems on Davis manifolds with finite 
nerve, the beta function has the qualitative form:
\begin{equation}\label{eq:beta_form}
    \beta(\Gamma) = \frac{1}{\nu} \cdot \Gamma \cdot 
    (\Gamma - 1) \cdot h(\Gamma)
\end{equation}
where:
\begin{itemize}
    \item $\nu$ is the correlation length exponent of the completion 
    transition (Branch~VI, S7);
    \item $h : (0, \infty) \to (0, \infty)$ is a smooth positive 
    function with $h(1) = 1$, depending on the universality class;
    \item the factor $\Gamma$ ensures $\beta(0) = 0$ 
    (underdetermined fixed point);
    \item the factor $(\Gamma - 1)$ ensures $\beta(1) = 0$ 
    (critical fixed point).
\end{itemize}
\end{conjecture}

\begin{remark}[Motivation for the Beta Function 
Form]\label{rem:beta_motivation}
The qualitative form~\eqref{eq:beta_form} is not derived from 
first principles but is constrained by three independent requirements:

\begin{enumerate}
    \item \textbf{Fixed-point structure.} The three regimes of the 
    Davis trichotomy (determined, critical, underdetermined) must be 
    fixed points of the RG flow. This requires $\beta(0) = \beta(1) 
    = 0$ and $\beta(\Gamma) \to +\infty$ as $\Gamma \to \infty$ 
    (so that $\Gamma = \infty$ is an attracting fixed point).
    
    \item \textbf{Flow direction.} Coarse-graining a determined 
    system ($\Gamma > 1$) averages out compatible sub-scale 
    solutions, producing an even-more-determined effective problem: 
    $\beta > 0$ for $\Gamma > 1$. Coarse-graining an underdetermined 
    system ($\Gamma < 1$) reveals even more freedom at larger scales: 
    $\beta < 0$ for $0 < \Gamma < 1$.
    
    \item \textbf{Consistency with Branch~VI.} At the critical point 
    $\Gamma = 1$, the derivative $\beta'(1) = 1/\nu$ must equal the 
    thermal exponent (the inverse correlation length exponent) to 
    recover the Branch~VI scaling $\xi \sim |\Gamma - 1|^{-\nu}$.
\end{enumerate}
The simplest function satisfying all three constraints is 
$(1/\nu) \cdot \Gamma \cdot (\Gamma - 1)$. The factor $h(\Gamma)$ 
absorbs higher-order corrections whose precise form depends on the 
microscopic constraint structure.
\end{remark}

The flow has the following properties:
\begin{enumerate}[label=(\alph*)]
    \item $\beta(\Gamma) < 0$ for $0 < \Gamma < 1$: the trichotomy 
    parameter \emph{decreases} under coarse-graining (flows toward 
    the underdetermined fixed point $\Gamma^* = 0$).
    
    \item $\beta(\Gamma) > 0$ for $\Gamma > 1$: $\Gamma$ 
    \emph{increases} under coarse-graining (flows toward the 
    determined fixed point $\Gamma^* = \infty$).
    
    \item $\beta(1) = 0$: the critical point is scale-invariant 
    (the constraint-to-freedom ratio does not change under 
    coarse-graining).
\end{enumerate}

\textbf{Derived from:} RG2 (semigroup) + RG3 (Laplacian) + 
Branch~VI, S7 (correlation length exponent $\nu$).

\begin{remark}[Discrete vs.\ Continuous Beta 
Function]\label{rem:discrete_beta}
The scale $\ell$ takes discrete values $\{0, 1, \ldots, L\}$ 
corresponding to the filtration levels of the partition 
(Definition~\ref{def:rg_semigroup}). The derivative 
$d\Gamma / d(\ln \ell)$ should therefore be understood as a 
discrete difference quotient:
\begin{equation}\label{eq:discrete_beta}
    \beta(\Gamma(\ell)) \approx 
    \frac{\Gamma(\ell + 1) - \Gamma(\ell)}
    {\ln(\ell + 1) - \ln(\ell)}
\end{equation}
The continuous form~\eqref{eq:beta_form} is a smooth interpolation 
through these discrete values. All downstream results (RG5 
fixed-point analysis, RG6 critical exponents, RG7 phase diagram) 
use only the sign structure and fixed-point values of $\beta$, 
which are well-defined in the discrete setting. In particular, 
$\beta(\Gamma^*) = 0$ means $\Gamma(\ell + 1) = \Gamma(\ell)$ at 
the fixed point, and the stability condition $\beta'(\Gamma^*) > 0$ 
means $|\Gamma(\ell + 1) - \Gamma^*| > |\Gamma(\ell) - \Gamma^*|$ 
(the deviation grows at each step). These are discrete dynamical 
statements requiring no differentiability.
\end{remark}

\textbf{Interpretation:} The beta function encodes how the 
``difficulty'' of a completion problem changes when viewed at 
different resolutions. Easy problems (determined regime) look 
progressively easier at coarser scales. Hard problems 
(underdetermined regime) look progressively harder. Critical 
problems look the same at every scale.

% ── RG5 ──
\section{RG5: Fixed Points and Stability}

\begin{conjecture}[Three Fixed Points of the Completion 
RG]\label{conj:fixed_points}
The beta function $\beta(\Gamma)$ of Conjecture~\ref{conj:beta} has 
exactly three fixed points in $[0, \infty]$:

\begin{center}
\renewcommand{\arraystretch}{1.4}
\begin{tabular}{lccll}
\textbf{Fixed point} & $\Gamma^*$ & \textbf{Stability} 
& \textbf{Physical regime} & \textbf{Cohomological signature} \\
\hline
Underdetermined & $0$ & IR stable & No effective constraints 
& $\dim \Hcech^0 \gg 1$, $\Hcech^1 = 0$ \\
Critical & $1$ & Unstable & Scale-invariant 
& spectral gap $= 0$ \\
Determined & $\infty$ & IR stable & Fully constrained 
& $\Hcech^0 = \{*\}$, $\Hcech^1 = 0$ \\
\end{tabular}
\end{center}
\end{conjecture}

\begin{remark}[The Determined Fixed Point at 
$\Gamma = \infty$]\label{rem:compactification}
The fixed point $\Gamma^* = \infty$ is a fixed point of the 
\emph{one-point compactification} $[0, \infty]$: writing 
$u = 1/\Gamma$, the beta function in the $u$-coordinate is 
$\beta_u = du / d(\ln \ell) = -u^2 \beta(1/u)$, which has a 
regular fixed point at $u^* = 0$ with stability exponent 
$d\beta_u/du|_{u=0} = -1/\nu < 0$ (stable, since $\beta_u < 0$ 
for small $u > 0$). The analysis of the determined regime is 
therefore well-defined in the compactified coordinate; the use 
of ``$\Gamma = \infty$'' is a notational convenience.
\end{remark}

\textbf{Stability analysis at $\Gamma^* = 1$.} Linearize 
$\beta(\Gamma) \approx (1/\nu) \cdot (\Gamma - 1) + 
O((\Gamma - 1)^2)$ near $\Gamma = 1$. The \emph{relevant 
eigenvalue} is:
\begin{equation}\label{eq:thermal_exponent}
    y_t = \left.\frac{d\beta}{d\Gamma}\right|_{\Gamma = 1} 
    = \frac{1}{\nu}
\end{equation}
This is the \emph{thermal exponent}. Since $y_t > 0$ ($\nu > 0$), 
the critical fixed point is \emph{unstable}: any perturbation 
$\Gamma = 1 + \varepsilon$ flows away exponentially under 
coarse-graining.

\textbf{Solution near criticality.} In the discrete setting 
(Remark~\ref{rem:discrete_beta}), the linearized RG map near 
$\Gamma^* = 1$ has a single relevant eigenvalue $\mu > 1$: each 
coarse-graining step multiplies the deviation from criticality by 
$\mu$. Writing $\delta\Gamma(\ell) = \Gamma(\ell) - 1$:
\begin{equation}\label{eq:discrete_recursion}
    \delta\Gamma(\ell + 1) = \mu \cdot \delta\Gamma(\ell) 
    + O(\delta\Gamma(\ell)^2)
\end{equation}
After $\ell$ discrete RG steps starting from scale $0$:
\begin{equation}\label{eq:flow_near_crit}
    \delta\Gamma(\ell) = \mu^\ell \cdot \delta\Gamma(0) 
    + O(\delta\Gamma(0)^2)
\end{equation}
This is purely discrete and requires no continuum limit.

\textbf{Defining $\nu$ from $\mu$.} The discrete eigenvalue $\mu$ is 
the \emph{primary observable} of the linearized RG map (in principle 
computable for any finite nerve). The correlation length exponent 
$\nu$ is \emph{defined} from $\mu$:
\begin{equation}\label{eq:nu_from_mu}
    \nu = \frac{1}{\ln \mu}
\end{equation}
so that $\mu = e^{1/\nu}$. The thermal exponent 
$y_t$~\eqref{eq:thermal_exponent} then satisfies $y_t = \ln \mu = 
1/\nu$. The continuous $\beta$-function (Remark~\ref{rem:discrete_beta}) 
reproduces this: $\beta'(1) = 1/\nu = \ln \mu$. There is no 
circularity: $\mu$ is computed from the discrete map, $\nu$ is derived, 
and the $\beta$-function's slope at the fixed point is a 
\emph{consequence}, not an input.

The continuous power-law form $\delta\Gamma \propto \ell^{1/\nu}$ 
arises only when one identifies the discrete step index $\ell$ 
with a continuous scale variable and writes $\mu^\ell = 
e^{\ell \ln \mu} = e^{(\ell/\nu) \cdot \nu \ln \mu}$. For 
logarithmically-spaced scales, $\mu^\ell = \ell^{1/\nu}$. 
This identification is a convenience, not a requirement: the 
discrete recursion~\eqref{eq:discrete_recursion} is the 
fundamental object.

\textbf{Correlation length from RG.} Define the correlation length 
$\xi$ as the scale $\ell^*$ at which 
$|\Gamma(\ell^*) - 1| \sim O(1)$:
\begin{equation}\label{eq:xi_from_rg}
    \xi \sim |\Gamma_0 - 1|^{-\nu}
\end{equation}
This recovers the Branch~VI result (S7) from the RG fixed-point 
structure.

\textbf{Derived from:} RG4 (beta function) + Branch~VI, S7 
(correlation length exponent $\nu$).

\textbf{Interpretation:} The three fixed points of the RG flow 
are the three regimes of the Davis trichotomy. The flow structure 
explains \emph{why} the trichotomy exists: it is the fixed-point 
structure of the coarse-graining operation. The instability of 
the critical point explains \emph{why} the phase transition is 
sharp: any perturbation away from $\Gamma = 1$ is amplified by 
the RG flow.

% ── RG6 ──
\section{RG6: Critical Exponents from RG Eigenvalues}

\begin{conjecture}[Critical Exponents as RG 
Eigenvalues]\label{conj:rg_exponents}
All critical exponents of the Davis completion transition 
(Branch~VI) are determined by the eigenvalues of the linearized 
RG transformation at the critical fixed point $\Gamma^* = 1$.

The completion transition has one relevant direction with 
eigenvalue:
\begin{equation}\label{eq:yt}
    y_t = \frac{1}{\nu}
\end{equation}
(the thermal exponent, controlling the distance $\Gamma - 1$ 
from criticality). All other directions are irrelevant 
(eigenvalues $< 0$), representing sub-scale constraint 
details that do not affect the critical behavior.

The standard scaling relations (derived, not independent) 
give all Branch~VI exponents from $y_t$ and the effective 
dimension $d_H$ of the constraint graph:
\begin{align}
    \alpha &= 2 - \nu \cdot d_H 
    && \text{(specific heat, Branch~VI S4)}
    \label{eq:hyperscaling}\\
    \gamma &= \nu(2 - \eta) 
    && \text{(susceptibility, Fisher's relation)}
    \label{eq:fisher}\\
    \beta_{\mathrm{mag}} &= \frac{\nu(d_H - 2 + \eta)}{2} 
    && \text{(order parameter)}
    \label{eq:beta_mag}
\end{align}
where $\eta$ is the anomalous dimension (correction to mean-field 
scaling of the \v{C}ech cochain two-point function).
\end{conjecture}

\begin{remark}[Upper Critical Dimension]\label{rem:upper_critical}
The hyperscaling relation~\eqref{eq:hyperscaling} holds when 
$d_H < d_{\mathrm{uc}}$, where $d_{\mathrm{uc}}$ is the 
\emph{upper critical dimension} of the completion transition. For 
$d_H \geq d_{\mathrm{uc}}$, mean-field exponents apply: $\alpha = 
0$, $\nu = 1/2$, $\eta = 0$, with possible logarithmic corrections 
at $d_H = d_{\mathrm{uc}}$. Completion problems on random graphs 
(random $k$-SAT, Erd\H{o}s--R\'enyi CSPs) have $d_H \to \infty$ 
and lie in the mean-field universality class. The value of 
$d_{\mathrm{uc}}$ for the Davis completion transition is an open 
question; by analogy with the Ising model, we conjecture 
$d_{\mathrm{uc}} = 4$, but the correct value may differ due to the 
non-standard interaction structure of constraint graphs.
\end{remark}

\textbf{Key prediction.} The anomalous dimension $\eta$ controls 
the decay of the \v{C}ech cochain correlator:
\begin{equation}\label{eq:anomalous_decay}
    G(v, w) = \langle \sigma(v) \, \sigma(w) \rangle_\beta 
    \sim d_\nerve(v, w)^{-(d_H - 2 + \eta)}
\end{equation}
at the critical point, where $d_\nerve$ is the graph metric on the 
nerve. Mean-field theory gives $\eta = 0$; nonzero $\eta$ arises 
from nonlinear interactions between constraint patches (the 
Baker--Campbell--Hausdorff corrections of Branch~VII, SH5).

\textbf{Derived from:} RG5 (fixed-point structure) + Branch~VI, S4 
(specific heat exponent $\alpha$), S7 (correlation length exponent 
$\nu$) + standard scaling relations.

\textbf{Interpretation:} The critical exponents are not free 
parameters --- they are \emph{eigenvalues} of the linearized RG at 
the critical fixed point. This is the sense in which the phase 
transition is ``universal'': systems with the same RG fixed point 
share the same eigenvalues and therefore the same critical behavior.

\begin{remark}[One Relevant Direction 
Hypothesis]\label{rem:one_relevant}
The conjecture that $\Gamma$ captures the \emph{only} relevant 
direction near the critical fixed point is the universality 
hypothesis. Irrelevant directions (sub-scale constraint details, 
gauge-equivalent configurations) flow to zero under 
coarse-graining and do not affect the critical behavior. This is 
the standard assumption in Wilson's RG and is supported by the 
observation that microscopically different CSPs (Sudoku, $k$-SAT, 
graph coloring) exhibit qualitatively similar phase transitions 
near their respective critical points.
\end{remark}

% ── RG7 ──
\section{RG7: The Phase Diagram as RG Flow Diagram}

\begin{conjecture}[Phase Diagram Flow 
Lines]\label{conj:phase_flow}
The phase diagram of the parent framework (Conjecture~50.1) in the 
$(\rho, \Khat)$ plane acquires RG flow lines. Each initial 
condition $(\rho_0, \Khat_0)$ flows under $T_\ell$ to a 
trajectory $(\rho(\ell), \Khat(\ell))$ with:

\begin{enumerate}[label=(\alph*)]
    \item \textbf{Above the critical line} ($\Gamma > 1$): flow 
    lines converge to the determined fixed point ($\rho \to \infty$, 
    $\Khat \to 0$) --- constraints dominate, effective curvature 
    vanishes.
    
    \item \textbf{Below the critical line} ($\Gamma < 1$): flow 
    lines converge to the underdetermined fixed point 
    ($\rho \to 0$, $\Khat \to 0$) --- both constraints and 
    curvature vanish (free theory).
    
    \item \textbf{On the critical line} ($\Gamma = 1$): flow lines 
    converge to the critical fixed point --- a unique point 
    $(\rho^*, \Khat^*)$ where the constraint-to-freedom ratio is 
    scale-invariant.
\end{enumerate}

The critical line $\rho_{\mathrm{crit}}(\Khat)$ is the 
\emph{separatrix} (co-dimension~1 invariant manifold) of the RG 
flow: the set of initial conditions flowing to the critical fixed 
point. Since the critical fixed point has exactly one relevant 
direction ($y_t = 1/\nu$), the separatrix is co-dimension~1, 
meaning a single parameter ($\Gamma$) must be tuned to reach 
criticality.
\end{conjecture}

\begin{remark}[RG-Improved Critical Line]\label{rem:rg_improved}
The RG flow provides corrections to the critical line at finite 
system size. Let $\omega > 0$ be the leading irrelevant exponent 
(the magnitude of the largest negative eigenvalue of the linearized 
RG). Then:
\begin{equation}\label{eq:rg_improved_line}
    \rho_{\mathrm{crit}}(\Khat, \ell) = 
    \rho_{\mathrm{crit}}^{(0)}(\Khat) + c \cdot \ell^{-\omega} 
    + \cdots
\end{equation}
At $\ell \to \infty$ (thermodynamic limit), corrections vanish and 
the critical line is sharp. At finite $\ell$ (finite problem size), 
the critical line is broadened by an amount $\sim \ell^{-\omega}$.
\end{remark}

\textbf{Derived from:} RG4 (beta function) + RG5 (fixed-point 
structure) + Parent Conjecture~50.1 (phase diagram).

\textbf{Interpretation:} The phase diagram is the \emph{projection} 
of the RG flow onto the $(\rho, \Khat)$ plane. The critical line 
is not an arbitrary boundary but the separatrix of a dynamical 
system --- the unique set of initial conditions balanced between 
the two stable fixed points.

% ── RG8 ──
\section{RG8: Monotonicity (C-Theorem Analog)}

\begin{theorem}[Completion Entropy 
Monotonicity]\label{thm:c_theorem}
Define the \emph{completion entropy} at scale $\ell$:
\begin{equation}\label{eq:completion_entropy}
    C(\ell) = \ln Z_0(\beta, \ell) 
    = \ln \tr \exp(-\beta \, \Delta_0^{(\ell)})
\end{equation}
where $Z_0(\beta, \ell)$ is the degree-0 partition function 
(Branch~VII, SH7) evaluated on the coarsened \v{C}ech Laplacian 
$\Delta_0^{(\ell)}$ with the inherited inner product 
(Definition~\ref{def:coarsened_cech}).

Then $C(\ell)$ is monotonically non-increasing under 
coarse-graining:
\begin{equation}\label{eq:c_monotone}
    C(\ell_2) \leq C(\ell_1) \qquad \text{for } \ell_2 > \ell_1
\end{equation}
\end{theorem}

\begin{proof}
We first establish a self-contained lemma, then apply it.

\begin{lemma}[Trace Monotonicity Under Subspace 
Restriction]\label{lem:trace_mono}
Let $A$ be an $n \times n$ positive semidefinite Hermitian matrix 
with eigenvalues $\lambda_1 \leq \cdots \leq \lambda_n$. Let 
$V \subseteq \mathbb{R}^n$ be a subspace of dimension $m \leq n$, 
and let $B = P^\dagger A P$ be the $m \times m$ compression, where 
$P : V \hookrightarrow \mathbb{R}^n$ is the inclusion. Let 
$\mu_1 \leq \cdots \leq \mu_m$ be the eigenvalues of $B$. Then 
for all $\beta > 0$:
\[
    \sum_{j=1}^{m} e^{-\beta \mu_j} 
    \leq \sum_{j=1}^{n} e^{-\beta \lambda_j}
\]
\end{lemma}

\begin{proof}[Proof of Lemma]
\textbf{Step 1} (Cauchy--Poincar\'e interlace theorem). For a 
Hermitian matrix $A$ with eigenvalues 
$\lambda_1 \leq \cdots \leq \lambda_n$ and its compression $B$ 
to an $m$-dimensional subspace with eigenvalues 
$\mu_1 \leq \cdots \leq \mu_m$, we have:
\begin{equation}\label{eq:interlace_explicit}
    \lambda_j \leq \mu_j \leq \lambda_{j + n - m} 
    \qquad \text{for } j = 1, \ldots, m
\end{equation}
In particular, $\mu_j \geq \lambda_j$ for all $j$. This is 
standard and proved via the min-max characterization of eigenvalues. 
Eigenvalues are counted with multiplicity throughout.

\textbf{Step 2} (term-by-term comparison). Since $\mu_j \geq 
\lambda_j$ and the function $x \mapsto e^{-\beta x}$ is strictly 
decreasing for $\beta > 0$:
\[
    e^{-\beta \mu_j} \leq e^{-\beta \lambda_j} 
    \qquad \text{for } j = 1, \ldots, m
\]
Summing over $j = 1, \ldots, m$:
\begin{equation}\label{eq:step2}
    \sum_{j=1}^{m} e^{-\beta \mu_j} 
    \leq \sum_{j=1}^{m} e^{-\beta \lambda_j}
\end{equation}

\textbf{Step 3} (dimension counting). Since $e^{-\beta \lambda_j} 
> 0$ for all $j$ and $m \leq n$:
\begin{equation}\label{eq:step3}
    \sum_{j=1}^{m} e^{-\beta \lambda_j} 
    \leq \sum_{j=1}^{n} e^{-\beta \lambda_j}
\end{equation}

Combining~\eqref{eq:step2} and~\eqref{eq:step3} gives the result.
No operator monotonicity, L\"owner order, majorization, or 
functional calculus is used --- only the Cauchy--Poincar\'e 
interlace theorem (a statement about finite sequences of real 
numbers) and the monotonicity of the real-valued function 
$e^{-\beta x}$.
\end{proof}

\textbf{Application to completion entropy.} By 
Definition~\ref{def:coarsened_cech}, the coarsened Laplacian is 
$\Delta_0^{(\ell)} = P_\ell^{(0)\dagger} \Delta_0 P_\ell^{(0)}$, 
the compression of $\Delta_0$ to the subspace $V_\ell$ of 
block-compatible 0-cochains. Since $\partition_{\ell_1} 
\preccurlyeq \partition_{\ell_2}$, we have $V_{\ell_2} \subseteq 
V_{\ell_1}$. The \v{C}ech Laplacian is positive semidefinite 
($\Delta_0 = (\delta^0)^\dagger \delta^0$, a sum of squares). 
Applying Lemma~\ref{lem:trace_mono} with $A = \Delta_0^{(\ell_1)}$ 
restricted to $V_{\ell_2} \subseteq V_{\ell_1}$ gives:
\[
    \tr \exp(-\beta \, \Delta_0^{(\ell_2)}) 
    \leq \tr \exp(-\beta \, \Delta_0^{(\ell_1)})
\]
Taking logarithms: $C(\ell_2) \leq C(\ell_1)$.

\textbf{Note (Remark~\ref{rem:compression_vs_coarse}).} The 
completion entropy uses the \emph{compression} $P_\ell^\dagger 
\Delta_0 P_\ell$, not the bare graph Laplacian of $\nerve_\ell$. 
These are different operators with different spectra; the 
C-theorem holds for the compression because it preserves the 
nested-subspace structure required by Cauchy interlace.
\end{proof}

\begin{remark}[Physical Interpretation]\label{rem:c_physical}
Coarse-graining can only \emph{destroy} completion possibilities, 
never create new ones (Theorem~\ref{thm:rg_inject}: the restriction 
$r_\ell$ is injective but not surjective). The completion entropy 
$C(\ell)$ counts (in a thermal-weighted sense) the effective number 
of completions visible at scale $\ell$. This number decreases 
monotonically:

\begin{itemize}
    \item At the determined fixed point: $C \to 0$ (one solution, 
    $Z_0 \to 1$).
    \item At the underdetermined fixed point: $C$ is maximal (many 
    solutions, large $Z_0$).
    \item At the critical fixed point: $C = C_{\mathrm{crit}}$, a 
    \emph{universal constant} that depends only on the universality 
    class. This universality of $C_{\mathrm{crit}}$ is conjectural 
    and analogous to Zamolodchikov's c-theorem for 2D conformal 
    field theories.
\end{itemize}
\end{remark}

\textbf{Derived from:} RG2 (semigroup, nested filtration) + 
Branch~VII, SH7 (partition function on \v{C}ech complex) + Cauchy 
interlace theorem.

\textbf{Interpretation:} The C-theorem provides an ``arrow of 
time'' for the RG flow. Information about completions flows 
from fine to coarse scales and is irreversibly lost. The 
monotonicity of $C(\ell)$ is the information-theoretic content of 
coarse-graining: you cannot recover fine-scale completion data from 
coarse-scale observations.

% ═══════════════════════════════════════════════════════
\newpage
\part{First-Order Derivations (RG1a--RG7a)}
% ═══════════════════════════════════════════════════════

\noindent
\textbf{Linearization convention.} All results in Parts~II and~III 
operate on the linearized presheaf $\sheaf^{\mathrm{lin}}$ 
(Branch~VII, abelian regime). The derived pushforward, spectral 
sequences, anomalous dimensions, and partition function computations 
require the additive coboundary structure. The set-valued 
coarse-graining of Part~I (RG1--RG2) does not; only the spectral 
analysis (RG3 onward) requires linearization.

% ── RG1a ──
\section{RG1a: Derived Pushforward and RG Spectral Sequence}

\begin{theorem}[Leray Spectral Sequence of 
Coarse-Graining]\label{thm:leray_rg}
The coarse-graining map $\pi_\ell : \nerve \to \nerve_\ell$ induces 
a derived pushforward $R\pi_{\ell *} : D^+(\nerve) \to 
D^+(\nerve_\ell)$. The higher direct images $R^q \pi_{\ell *} 
\sheaf$ form presheaves on $\nerve_\ell$ with stalks:
\begin{equation}\label{eq:higher_direct}
    (R^q \pi_{\ell *} \sheaf)(\block_a) 
    = H^q\!\left(\pi_\ell^{-1}(\block_a),\; 
    \sheaf\big|_{\pi_\ell^{-1}(\block_a)}\right)
\end{equation}
measuring the degree-$q$ cohomological obstruction within each 
block. There is a Leray spectral sequence:
\begin{equation}\label{eq:leray}
    E_2^{p,q} = \Hcech^p(\nerve_\ell,\; R^q \pi_{\ell *} \sheaf) 
    \;\Longrightarrow\; \Hcech^{p+q}(\nerve, \sheaf)
\end{equation}
\end{theorem}

\begin{proof}
This is the standard Leray spectral sequence applied to the 
simplicial map $\pi_\ell$ and the presheaf $\sheaf$. The 
construction requires only that $\pi_\ell$ is a surjective 
simplicial map between finite simplicial complexes and that 
$\sheaf$ is a presheaf of abelian groups (which holds for 
$\sheaf^{\mathrm{lin}}$). The spectral sequence converges because 
$\nerve$ and $\nerve_\ell$ are finite-dimensional.
\end{proof}

\begin{conjecture}[RG Interpretation of the Spectral 
Sequence]\label{conj:rg_spectral}
The terms of the Leray spectral sequence have the following RG 
interpretation:

\begin{enumerate}[label=(\alph*)]
    \item $R^0 \pi_{\ell *} \sheaf = \sheaf_\ell$: the 
    \emph{renormalized presheaf} (effective theory at scale $\ell$).
    
    \item $R^q \pi_{\ell *} \sheaf$ for $q \geq 1$: the 
    \emph{integrated-out degrees of freedom}. These are sub-scale 
    obstructions absorbed into the effective theory.
    
    \item \textbf{Leray condition at scale $\ell$:} If 
    $R^q \pi_{\ell *} \sheaf = 0$ for all $q \geq 1$ (every block 
    is internally acyclic), the spectral sequence degenerates at 
    $E_2$ and:
    \begin{equation}\label{eq:leray_degen}
        \Hcech^n(\nerve, \sheaf) \cong 
        \Hcech^n(\nerve_\ell, \sheaf_\ell) \qquad \text{for all } n
    \end{equation}
    No cohomological information is lost by coarse-graining.
\end{enumerate}

\textbf{RG relevance classification.} An obstruction class 
$[\alpha] \in \Hcech^1(\nerve, \sheaf)$ is:
\begin{itemize}
    \item \emph{Relevant} if its image in 
    $\Hcech^1(\nerve_\ell, \sheaf_\ell)$ is nonzero for all $\ell$ 
    (survives all coarse-graining).
    \item \emph{Irrelevant} if its image vanishes at some finite 
    $\ell$ (a sub-scale obstruction).
    \item \emph{Marginal} if it maps to zero in $\Hcech^1$ but 
    appears in $\Hcech^0(\nerve_\ell, R^1 \pi_{\ell *} \sheaf)$ 
    (visible only through the spectral sequence correction).
\end{itemize}
\end{conjecture}

\textbf{Derived from:} RG1 (coarse-graining) + Branch~VII, SH5a 
(\v{C}ech-to-derived spectral sequence) + X10 (renormalization 
as derived pushforward).

% ── RG2a ──
\section{RG2a: Anomalous Dimensions and Cocycle Scaling}

\begin{conjecture}[Anomalous Scaling at 
Criticality]\label{conj:anomalous}
At the critical fixed point $\Gamma^* = 1$, the \v{C}ech cochains 
acquire \emph{anomalous scaling dimensions}. The two-point 
function on 0-cochains is the Green's function of the \v{C}ech 
Laplacian:
\begin{equation}\label{eq:green}
    G(v, w) = \langle \sigma(v) \, \sigma(w) \rangle_\beta 
    = (\Delta_0^{-1})(v, w)
\end{equation}
At the critical point, $G(v, w)$ decays as a power law:
\begin{equation}\label{eq:anomalous}
    G(v, w) \sim d_\nerve(v, w)^{-(d_H - 2 + \eta)}
\end{equation}
where $d_H$ is the effective dimension and $\eta$ is the anomalous 
dimension.

The scaling dimension of a $p$-cochain field is:
\begin{equation}\label{eq:cochain_dimension}
    [\sigma_p] = \frac{d_H - 2 + \eta}{2} + p
\end{equation}
Mean-field theory gives $\eta = 0$. Nonzero $\eta$ arises from 
the BCH quadratic corrections (Branch~VII, SH5):
\begin{equation}\label{eq:eta_bch}
    \eta = c \cdot \|\alpha_c\|^2 + O(\|\alpha_c\|^4)
\end{equation}
where $\|\alpha_c\|$ is the typical cocycle norm at criticality 
and $c$ is a geometry-dependent constant.
\end{conjecture}

\textbf{Derived from:} RG6 (critical exponents) + Branch~VII, SH5 
(BCH control).

% ── RG3a ──
\section{RG3a: Crossover Scaling Functions}

\begin{conjecture}[Universal Scaling 
Form]\label{conj:crossover}
The free energy density $f = F/n$ obeys a universal scaling form 
near the critical point $\Gamma^* = 1$:
\begin{equation}\label{eq:scaling_form}
    f(\Gamma, \beta) = |\Gamma - 1|^{2 - \alpha} \cdot 
    \Phi_\pm\!\left(\beta \cdot |\Gamma - 1|^{\nu \cdot d_H}\right)
\end{equation}
where $\Phi_\pm$ are universal scaling functions (one for 
$\Gamma > 1$, one for $\Gamma < 1$), $\alpha = 2 - \nu \cdot d_H$ 
(hyperscaling), and $\beta$ is the inverse temperature of the 
partition function (Branch~VII, SH7).

Three regimes emerge:
\begin{enumerate}[label=(\alph*)]
    \item $\beta |\Gamma - 1|^{\nu d_H} \gg 1$ (ordered): deep in 
    a phase, mean-field behavior, exponential correlation decay.
    
    \item $\beta |\Gamma - 1|^{\nu d_H} \ll 1$ (critical): 
    fluctuation-dominated, power-law correlations, anomalous 
    exponents.
    
    \item $\beta |\Gamma - 1|^{\nu d_H} \sim 1$ (crossover): 
    $\Phi_\pm$ interpolates smoothly between ordered and critical 
    behavior.
\end{enumerate}
\end{conjecture}

\textbf{Derived from:} RG5 (fixed-point structure) + RG6 (critical 
exponents) + Branch~VI (thermodynamic formalism).

% ── RG4a ──
\section{RG4a: Corrections to Scaling}

\begin{conjecture}[Irrelevant Operators and Correction 
Exponents]\label{conj:corrections}
The leading \emph{correction-to-scaling exponent} $\omega > 0$ is 
the magnitude of the largest irrelevant eigenvalue of the 
linearized RG at $\Gamma^* = 1$. Observables receive power-law 
corrections:
\begin{align}
    \xi &= \xi_0 \, |\Gamma - 1|^{-\nu} 
    \left(1 + c_1 \, |\Gamma - 1|^{\omega \nu} + \cdots\right)
    \label{eq:xi_correction}\\
    C_V &= C_0 \, |\Gamma - 1|^{-\alpha} 
    \left(1 + c_2 \, |\Gamma - 1|^{\omega \nu} + \cdots\right)
    \label{eq:cv_correction}
\end{align}

Irrelevant operators in the Davis context are sub-scale constraint 
arrangements that do not affect long-distance completion physics: 
cell permutations within a Sudoku box, variable relabeling in 
$k$-SAT, local gauge transformations of the cocycle. The 
correction exponent can be computed from the \v{C}ech Laplacian 
spectrum at the critical point:
\begin{equation}\label{eq:omega_spectral}
    \omega = \frac{\ln(\lambda_2 / \lambda_1)}{\ln \ell}
\end{equation}
where $\lambda_1, \lambda_2$ are the two smallest nonzero 
eigenvalues of $\Delta_1^{(\ell)}$ at criticality.
\end{conjecture}

\textbf{Derived from:} RG5 (fixed-point structure) + RG6 (critical 
exponents).

% ── RG5a ──
\section{RG5a: Universality Classes of Completion Problems}

\begin{conjecture}[Universality 
Classification]\label{conj:universality}
Completion problems flowing to the same critical fixed point under 
RG share the same critical exponents --- the same 
\emph{universality class}. The universality class is determined by:
\begin{enumerate}[label=(\roman*)]
    \item $d_H$: the effective dimension of the constraint graph 
    (Hausdorff dimension of the nerve for structured problems, or 
    $d_H \to \infty$ for random graphs).
    
    \item The \emph{symmetry group} of the constraint structure 
    (automorphism group of the nerve acting on the presheaf).
    
    \item The \emph{range of interaction} (whether constraints have 
    finite-range or long-range overlap structure).
\end{enumerate}

Conjectured universality classes include:

\begin{center}
\renewcommand{\arraystretch}{1.3}
\begin{tabular}{llp{3.8cm}}
\textbf{Class} & $d_H$ & \textbf{Examples} \\
\hline
Short-range lattice & $2$ & Sudoku, Latin squares, crosswords \\
Random graph (mean-field) & $\to \infty$ & Random $k$-SAT, 
Erd\H{o}s--R\'enyi CSP \\
Hierarchical & $\log_b a$ & Nested constraint trees \\
Long-range & $d_{\mathrm{eff}} < d$ & Attention-connected 
transformer blocks \\
\end{tabular}
\end{center}
\end{conjecture}

\textbf{Derived from:} RG6 (critical exponents from RG eigenvalues).

\textbf{Interpretation:} Universality explains why microscopically 
different completion problems exhibit the same critical behavior: 
they flow to the same RG fixed point. The irrelevant operators 
(microscopic details) wash out under coarse-graining, and only the 
relevant direction ($\Gamma - 1$) and the fixed-point eigenvalue 
($1/\nu$) determine the critical behavior.

% ── RG6a ──
\section{RG6a: RG-Improved Partition Function}

\begin{conjecture}[RG Resummation of the Partition 
Function]\label{conj:rg_partition}
The bare partition function $Z_0(\beta)$ (Branch~VII, SH7) 
receives RG corrections. The \emph{RG-improved partition function} 
replaces bare parameters with running couplings:
\begin{equation}\label{eq:rg_partition}
    Z_{\mathrm{RG}}(\beta) = \tr \exp(-\beta \, 
    \Delta_0(g(\beta)))
\end{equation}
where $g(\beta) = \{m(\ell_T)^2, n(\ell_T), s(\ell_T), \ldots\}$ 
are the running couplings evaluated at the \emph{thermal scale} 
$\ell_T = \beta^{1/d_H}$. The running couplings satisfy the RG 
equations (Conjecture~\ref{conj:beta}) at $\ell = \ell_T$.

The RG-improved partition function:
\begin{enumerate}[label=(\alph*)]
    \item Captures the correct critical scaling (Branch~VI exponents 
    reproduced by construction).
    \item Agrees with the bare $Z_0$ at high temperature 
    ($\beta \to 0$, $\ell_T \to 0$, no running).
    \item Agrees with the ground-state count at low temperature 
    ($\beta \to \infty$, $\ell_T \to \infty$, fully coarsened).
\end{enumerate}
\end{conjecture}

\textbf{Derived from:} RG2 (semigroup) + RG4 (beta function) + 
Branch~VII, SH7 (partition function).

% ── RG7a ──
\section{RG7a: Euler Characteristic as RG Invariant}

\begin{theorem}[Topological Invariance of Euler 
Characteristic]\label{thm:euler_invariant}
Under the hypotheses of Branch~VII, SH7 (finite nerve, 
finite-dimensional stalks) and the \emph{Leray-acyclicity 
condition} $R^q \pi_{\ell *} \sheaf = 0$ for all $q \geq 1$ 
(every block $\block_a$ is internally acyclic: 
$H^q(\pi_\ell^{-1}(\block_a), \sheaf|_{\pi_\ell^{-1}(\block_a)}) 
= 0$ for $q \geq 1$), the sheaf Euler characteristic is invariant 
under coarse-graining:
\begin{equation}\label{eq:euler_invariant}
    \chi(\sheaf) = \chi(\sheaf_\ell) \qquad 
    \text{(under Leray-acyclicity)}
\end{equation}
\end{theorem}

\begin{proof}
The Leray spectral sequence 
(Theorem~\ref{thm:leray_rg},~\eqref{eq:leray}) has $E_2^{p,q} = 
\Hcech^p(\nerve_\ell, R^q \pi_{\ell *} \sheaf)$ converging to 
$\Hcech^{p+q}(\nerve, \sheaf)$. The Euler characteristic of each 
page of a spectral sequence is the same (the alternating sum of 
dimensions is preserved by the differentials). Therefore:
\begin{equation}\label{eq:euler_pages}
    \chi(\sheaf) = \sum_n (-1)^n \dim \Hcech^n(\nerve, \sheaf) 
    = \sum_{p,q} (-1)^{p+q} \dim E_2^{p,q}
\end{equation}
Expanding the $E_2$ page:
\begin{equation}\label{eq:euler_expand}
    \sum_{p,q} (-1)^{p+q} \dim E_2^{p,q} 
    = \sum_q (-1)^q \chi(R^q \pi_{\ell *} \sheaf)
\end{equation}
Under the Leray-acyclicity hypothesis, $R^q \pi_{\ell *} \sheaf = 0$ 
for all $q \geq 1$. Only the $q = 0$ term survives:
\[
    \chi(\sheaf) = \chi(R^0 \pi_{\ell *} \sheaf) 
    = \chi(\sheaf_\ell)
\]
since $R^0 \pi_{\ell *} \sheaf = \sheaf_\ell$ by construction 
(Definition~\ref{def:coarse_grain}).
\end{proof}

\begin{remark}[When Leray-Acyclicity Holds]
\label{rem:leray_acyclic}
The Leray-acyclicity condition requires that each block 
$\block_a$ is ``internally coherent'': any local sections on the 
constraints within $\block_a$ that agree on pairwise overlaps 
already agree on all higher overlaps. For blocks of small diameter 
(low $\ell$), this is physically reasonable: constraints in a 
small neighborhood interact simply. The condition can fail for 
large blocks that combine constraints with complex internal 
topology.

When the condition fails, the spectral sequence still gives 
$\chi(\sheaf) = \sum_q (-1)^q \chi(R^q \pi_{\ell *} \sheaf)$, 
but the individual terms with $q \geq 1$ need not vanish, and 
$\chi(\sheaf) \neq \chi(\sheaf_\ell)$ in general. The difference 
$\chi(\sheaf) - \chi(\sheaf_\ell) = \sum_{q \geq 1} (-1)^q 
\chi(R^q \pi_{\ell *} \sheaf)$ measures the sub-scale 
cohomological information invisible at scale $\ell$.
\end{remark}

\begin{remark}[Physical Interpretation]
\label{rem:euler_physical}
The Euler characteristic $\chi(\sheaf) = \sum (-1)^p \dim 
\Hcech^p(\nerve, \sheaf)$ counts the net degrees of freedom 
(solutions minus obstructions). Under the Leray-acyclicity 
condition of Theorem~\ref{thm:euler_invariant}, this is an RG 
invariant: $\chi(\sheaf) = \chi(\sheaf_\ell)$. The RG flow 
changes the distribution of information across cohomological 
degrees (solutions and obstructions can individually change) but 
the net balance is fixed.

When the Leray condition fails, the general identity 
$\chi(\sheaf) = \sum_q (-1)^q \chi(R^q \pi_{\ell *} \sheaf)$ 
still holds, but $\chi(\sheaf_\ell)$ alone does not capture the 
full Euler characteristic --- the higher direct images contribute.

The McKean--Singer formula in the Davis context states: the 
supertrace $\sum (-1)^p \tr(e^{-\beta \Delta_p})$ equals 
$\chi(\sheaf)$ independently of $\beta$. Under Leray-acyclicity, 
this is also independent of the coarse-graining scale $\ell$.
\end{remark}

\textbf{Derived from:} RG1a (Leray spectral sequence) + 
Branch~VII, SH7 (Euler characteristic).

% ═══════════════════════════════════════════════════════
\newpage
\part{Cross-Synthesis with Parent Framework and Branches~VI--VII}
% ═══════════════════════════════════════════════════════

% ── Y1 ──
\section{Y1: Transformer Layers as RG Steps}

\begin{conjecture}[Layer-Scale 
Correspondence]\label{conj:layers_rg}
In a transformer with $L$ layers processing a completion task, 
each layer $\ell \in \{1, \ldots, L\}$ implements one RG 
coarse-graining step $T_{\ell/L \cdot \ell_{\max}}$, where 
$\ell_{\max}$ is the maximum scale (set by the context window 
length). The residual stream at layer $\ell$ carries the 
scale-dependent trichotomy parameter $\Gamma(\ell)$.

This yields three layer-wise predictions:
\begin{enumerate}[label=(\alph*)]
    \item \textbf{Easy completions ($\Gamma_0 > 1$):} 
    $\Gamma(\ell)$ increases rapidly through layers. The residual 
    stream converges early and attention patterns become sparse 
    (few heads active in later layers). This matches empirical 
    observations that transformers ``solve'' easy tasks in early 
    layers.
    
    \item \textbf{Hard completions ($\Gamma_0 \approx 1$):} 
    $\Gamma(\ell)$ stays near $1$ through many layers. Attention 
    patterns exhibit long-range structure at all scales. This 
    matches the critical slowing-down observed in transformer 
    processing of ambiguous or complex inputs.
    
    \item \textbf{Impossible completions ($\Gamma_0 < 1$):} 
    $\Gamma(\ell)$ decreases through layers toward the 
    underdetermined fixed point. No convergence occurs, and the 
    model hallucinates (generates output from the underdetermined 
    regime where many completions exist but none is uniquely 
    correct).
\end{enumerate}
\end{conjecture}

\textbf{Derived from:} RG4 (beta function) + RG5 (fixed-point 
structure).

\textbf{Interpretation:} Hallucination is not a failure of the 
model per se --- it is the correct dynamical behavior of a system 
in the underdetermined basin of attraction: $\Gamma < 1$ flows 
toward $\Gamma^* = 0$, where all local completions are compatible 
and any output is ``valid'' at the coarsest scale. The model 
produces a completion because it must, but the completion is 
underdetermined by the input.

% ── Y2 ──
\section{Y2: Multi-Head Attention as Multi-Scale Parallel Transport}

\begin{conjecture}[Multi-Head Multi-Scale 
Decomposition]\label{conj:multihead}
Different attention heads within the same transformer layer operate 
at different effective scales, implementing parallel RG steps at 
multiple resolutions simultaneously:
\begin{equation}\label{eq:multihead_scale}
    \text{Head } h \text{ at layer } \ell 
    \quad\longleftrightarrow\quad T_{\ell_h}
\end{equation}
where $\ell_h = f(h, \ell)$ is a head-dependent effective scale.

The attention matrix $A^{(h)}$ of head $h$ has an effective range 
$\mathrm{range}(A^{(h)}) \sim \ell_h$: short-range heads (small 
$\ell_h$) handle local constraint compatibility, while long-range 
heads (large $\ell_h$) handle global coherence. The multi-head 
mechanism provides a \emph{multi-scale decomposition} of the RG 
flow within each layer.

\textbf{Prediction:} Transformers are more powerful than 
single-head architectures because they simultaneously process 
constraint information at multiple scales within each layer, rather 
than being limited to a single coarse-graining step.
\end{conjecture}

\textbf{Derived from:} Y1 (layer-scale correspondence) + 
Branch~VII, SH5 (parallel transport on the nerve).

% ── Y3 ──
\section{Y3: Scale-Dependent Completion Capacity}

\begin{conjecture}[Depth-Capacity 
Scaling]\label{conj:capacity}
The \emph{completion capacity} of a model at scale $\ell$ is:
\begin{equation}\label{eq:capacity}
    C_{\mathrm{cap}}(\ell) = \max\{\Gamma : \text{model can reach 
    determined fixed point starting from } \Gamma \text{ at scale } 
    \ell\}
\end{equation}
For a transformer with $L$ layers:
\begin{equation}\label{eq:capacity_depth}
    C_{\mathrm{cap}} \approx 1 + c \cdot L^{1/\nu}
\end{equation}
where $c$ is a model-dependent constant and $\nu$ is the 
correlation length exponent.

\textbf{Scaling prediction:} Doubling transformer depth increases 
capacity by a factor of $2^{1/\nu}$. If $\nu \approx 1$ 
(mean-field), doubling depth doubles capacity. If $\nu > 1$, the 
return on depth diminishes (sublinear in layers).
\end{conjecture}

\textbf{Derived from:} Y1 (layers as RG steps) + RG5 (stability) 
+ RG6 (critical exponents).

% ── Y4 ──
\section{Y4: Residual Stream as RG Trajectory}

\begin{conjecture}[Residual Stream 
Dynamics]\label{conj:residual}
The residual stream $x_\ell$ at layer $\ell$ encodes the effective 
completion state at scale $\ell$. The layer-to-layer update:
\begin{equation}\label{eq:residual_update}
    x_{\ell+1} = x_\ell + \mathrm{Attn}(x_\ell) + 
    \mathrm{FFN}(x_\ell)
\end{equation}
implements an infinitesimal RG step:
\begin{equation}\label{eq:rg_step}
    (\nerve_{\ell+\delta\ell}, \sheaf_{\ell+\delta\ell}) 
    = T_{\delta\ell}(\nerve_\ell, \sheaf_\ell)
\end{equation}

The norm of the per-layer update 
$\|x_{\ell+1} - x_\ell\|$ is large when $\Gamma(\ell) \approx 1$ 
(critical regime, active processing) and small when 
$|\Gamma(\ell) - 1| \gg 0$ (deep in a phase, convergence). 

\textbf{Prediction:} Per-layer activation norms peak at layers 
where the effective completion problem is near criticality.
\end{conjecture}

\textbf{Derived from:} Y1 (layers as RG steps).

% ── Y5 ──
\section{Y5: Obstruction Persistence Under RG}

\begin{conjecture}[Persistence as RG 
Relevance]\label{conj:persistence}
The persistence diagram of the completion presheaf (Branch~VII, 
SH4b/X9) classifies obstruction classes by their \emph{RG 
relevance}. Each class $[\alpha] \in \Hcech^1(\nerve, \sheaf)$ has:
\begin{itemize}
    \item \emph{Birth scale} $\ell_{\mathrm{birth}}$: the coarsest 
    scale at which the obstruction appears.
    \item \emph{Death scale} $\ell_{\mathrm{death}}$: the finest 
    scale at which the obstruction can be resolved.
\end{itemize}

Persistent classes ($\ell_{\mathrm{death}} - \ell_{\mathrm{birth}}$ 
large) are \emph{relevant operators}: they survive coarse-graining 
and control long-distance completion physics. Ephemeral classes 
($\ell_{\mathrm{death}} - \ell_{\mathrm{birth}}$ small) are 
\emph{irrelevant operators}: sub-scale obstructions that do not 
affect the large-scale completion.

The difficulty spectrum $D(\tau)$ from Branch~VII, X9 acquires a 
scale-dependent form:
\begin{equation}\label{eq:difficulty_rg}
    D(\tau, \ell) = |\{[\alpha] : 
    \ell_{\mathrm{birth}}([\alpha]) \leq \ell 
    \leq \ell_{\mathrm{death}}([\alpha]),\; 
    \tau_{\mathrm{death}}([\alpha]) > \tau\}|
\end{equation}
The RG flow decreases $D(\tau, \ell)$ monotonically, consistent 
with the C-theorem (Theorem~\ref{thm:c_theorem}).
\end{conjecture}

\textbf{Derived from:} RG1a (relevance classification) + RG8 
(monotonicity) + Branch~VII, SH4b/X9 (persistence diagrams).

% ── Y6 ──
\section{Y6: Asymptotic Freedom}

\begin{conjecture}[Asymptotic Freedom of the Davis 
Theory]\label{conj:asymptotic_freedom}
The holonomy budget constraint (Parent framework, Branch~VII SH5(d)) 
acquires a scale-dependent form under RG:
\begin{equation}\label{eq:holonomy_budget_rg}
    \|\alpha\|_{\Cech^1}^{(\ell)} < \frac{\tau}{k(\ell)}
\end{equation}
where $k(\ell)$ is the maximum loop length in the coarsened nerve 
$\nerve_\ell$.

Coarse-graining reduces loop lengths under the following 
hypothesis:

\textbf{Loop contraction hypothesis.} For coarse-grainings whose 
blocks have bounded overlap (the induced partition of the 
1-skeleton is $\ell$-Lipschitz with respect to the graph metric), 
the maximum loop length satisfies:
\begin{equation}\label{eq:loop_contraction}
    k(\ell) \leq C \cdot k_0 / \ell
\end{equation}
for a constant $C$ depending on the local geometry of the nerve. 
This holds whenever the coarse-graining contracts graph distances 
by a factor proportional to $\ell$ (the block diameter), which is 
the case for structured nerves (lattice-like, hierarchical) but may 
fail for highly irregular graph topologies. In particular, 
expander-type nerves (where long geodesic cycles survive 
coarse-graining because high connectivity prevents cycle collapse) 
and graphs with girth $\gg$ diameter are excluded.

Under this hypothesis, the 
effective budget per cocycle \emph{increases} at larger 
scales:
\begin{equation}\label{eq:tau_eff}
    \tau_{\mathrm{eff}}(\ell) 
    = \frac{\tau \cdot \ell}{k_0}
\end{equation}

Moreover, in the linearized regime, the coarsened cocycle norm is 
bounded:
\begin{equation}\label{eq:cocycle_bound}
    \|\alpha_\ell(a, b)\| \leq \max\{\|\alpha(v, w)\| 
    : v \in \block_a,\; w \in \block_b,\; (v, w) \text{ edge}\} 
    \leq \|\alpha\|_{\Cech^1}
\end{equation}
since the block-level transition is determined by a single 
boundary overlap (or a finite composition of commuting transitions 
in the additive regime). The total accumulation around a coarsened 
loop of length $k(\ell)$ is thus bounded by:
\begin{equation}\label{eq:total_accumulation}
    k(\ell) \cdot \|\alpha_\ell\|_{\Cech^1} 
    \leq \frac{k_0}{\ell} \cdot \|\alpha\|_{\Cech^1} 
    = \frac{1}{\ell} \cdot k_0 \cdot \|\alpha\|_{\Cech^1}
\end{equation}
which shrinks with $\ell$.

\textbf{Conclusion:} At large scales, the effective coupling is 
weak and the linearized theory (Branch~VII, additive regime) 
becomes increasingly accurate. The BCH convergence condition 
$\|\alpha\| < (\log 2) / k(\ell)$ (Branch~VII, 
Remark~5.1) is automatically satisfied at sufficiently 
large $\ell$. The entire linearized apparatus of Branches~VI--VII 
is \emph{self-consistently justified in the infrared}.
\end{conjecture}

\textbf{Derived from:} Branch~VII, SH5 (holonomy, BCH control) + 
RG2 (scale-dependent parameters).

\textbf{Interpretation:} Asymptotic freedom means that the 
nonlinear corrections (BCH terms, group-valued cocycles) are 
\emph{ultraviolet effects}: they matter at fine scales where loop 
lengths are long and cocycle accumulation is large, but they wash 
out at coarse scales where loops are short and the linearized 
theory is valid. This is the inverse of the QCD pattern (where the 
coupling is weak at short distances and strong at long distances); 
in the Davis theory, the coupling is strong at short distances 
(fine-scale constraint interactions) and weak at long distances 
(coarse-scale effective theory).

% ── Y6a ──
\section{Y6a: Scale-Local Connections (Product Gauge Structure)}

\begin{conjecture}[Scale-Local 
Connections]\label{conj:scale_local}
A multi-scale transport architecture implementing $N$ successive 
RG steps (Y1) requires $N$ \emph{scale-local connections}, one per 
scale. That is, the connection (equivalently, the metric from which 
it derives) at scale $\ell$ must be the connection of the 
\emph{coarsened} geometry $\Delta_p^{(\ell)}$, not the original 
fine-scale geometry $\Delta_p$.

\textbf{Statement.} Let $(M, g)$ be a Davis manifold and 
$T_{\ell_1}, \ldots, T_{\ell_N}$ be $N$ RG coarse-graining 
steps at scales $\ell_1 < \cdots < \ell_N$. The parallel 
transport operator at scale $\ell_k$ is determined by the 
connection $\nabla^{(\ell_k)}$ derived from the compressed 
Laplacian $\Delta_p^{(\ell_k)} = P_{\ell_k}^\dagger \Delta_p 
P_{\ell_k}$ (RG3). This connection is generically distinct from 
$\nabla^{(\ell_j)}$ for $j \neq k$, because the compressed 
Laplacian has a different spectrum at each scale 
(Remark~\ref{rem:compression_vs_coarse}).

\textbf{Consequence.} A neural architecture implementing 
multi-scale transport cannot share a single learned metric across 
scales. Each transport block requires its own metric network 
$g^{(\ell_k)}$, which learns the effective geometry at scale 
$\ell_k$. The inter-scale coupling is provided by the 
coarse-graining map itself --- in a transformer, the attention 
layer between transport blocks.
\end{conjecture}

\textbf{Derived from:} Y1 (layers as RG steps) + RG3 (coarsened 
Laplacian, Remark~\ref{rem:compression_vs_coarse}: spectrum changes 
under coarsening) + Y6 (asymptotic freedom: coupling constants 
differ across scales).

\begin{remark}[Gauge-Theoretic 
Interpretation]\label{rem:product_gauge}
The scale-local connection structure has a natural interpretation 
in gauge theory. Each scale $\ell_k$ carries its own gauge group 
$G_k = O(d)$ (orthogonal transport in $d$ dimensions) with its own 
connection $\nabla^{(\ell_k)}$. The total gauge group of an 
$N$-scale architecture is the product:
\begin{equation}\label{eq:product_gauge}
    G = G_1 \times G_2 \times \cdots \times G_N
\end{equation}
with each factor possessing an independent connection (equivalently, 
an independent metric network). The factors do not interact 
directly: there is no mixed $G_i$--$G_j$ gauge field. Inter-scale 
coupling occurs through the \emph{matter fields} --- the hidden 
state representations $h_\ell$ that transform under the gauge 
group at their respective scale. The attention layer between 
scales $\ell_k$ and $\ell_{k+1}$ is the ``matter field'' that 
couples the two gauge sectors.

This structure extends naturally: adding a third scale adds 
$G_3$ to the product, with its own metric network and its own 
connection. The architectural complexity scales linearly with the 
number of scales, not combinatorially.
\end{remark}

\begin{remark}[Finite-Model 
Evidence]\label{rem:finite_model_y6a}
The Shidoku finite model (Appendix~\ref{app:shidoku_rg}) provides 
direct evidence. The 12-vertex nerve with its 4-block partition 
produces a compressed Laplacian with spectrum $\{0, 3, 6, 9\}$, 
while the bare $K_{2,2}$ Laplacian of the coarsened nerve has 
spectrum $\{0, 2, 2, 4\}$. These are genuinely different operators 
with different eigenvectors. Any transport process using the 
$\{0, 2, 2, 4\}$ geometry at the coarse scale would be solving 
the wrong variational problem --- the effective constraint 
structure at that scale is $\{0, 3, 6, 9\}$, not $\{0, 2, 2, 4\}$. 
A~fortiori, using the fine-scale spectrum (12~eigenvalues) at the 
coarse scale would be even more wrong: the coarsened geometry lives 
in a 4-dimensional subspace of the original 12-dimensional space, 
and the compressed Laplacian captures exactly the spectral content 
relevant to that subspace.
\end{remark}

% ── Y6b ──
\section{Y6b: Bounded Coarse-Scale Transport (Low-Rank Connection)}

The scale-local connection structure of Y6a introduces a new 
instability: the coarse manifold $\mathcal{M}_2$ receives 
\emph{contextual} base points (post-attention hidden states) rather 
than discrete vocabulary embeddings. We show that finite-difference 
Christoffel estimation diverges in this regime, and that a low-rank 
factored connection provides an architecturally bounded alternative.

\begin{proposition}[Divergence of FD Christoffels on Contextual 
Positions]\label{prop:fd_divergence}
Let $\mathcal{M}$ be a Davis manifold with learned metric 
$G(x) = L(x)L(x)^\top + \varepsilon I$ where $L(x)$ is a neural 
network. Define the finite-difference Christoffel symbols
\begin{equation}\label{eq:fd_christoffel}
    \hat{\Gamma}^k_{ij}(x) = \frac{1}{2} G^{kl}(x) \left(
    \frac{G_{li}(x + h e_j) - G_{li}(x - h e_j)}{2h}
    + \frac{G_{lj}(x + h e_i) - G_{lj}(x - h e_i)}{2h}
    - \frac{G_{ij}(x + h e_k) - G_{ij}(x - h e_k)}{2h}
    \right).
\end{equation}
When $x$ ranges over a continuous manifold of contextual hidden 
states $h_t \in \mathbb{R}^d$ (as opposed to a finite vocabulary 
$\{e_1, \ldots, e_V\}$), the FD estimates suffer a positive 
feedback loop:
\begin{enumerate}
    \item The metric network $L$ is cold-started on inputs $h_t$ 
    that are themselves a moving training target;
    \item Cold $L$ produces ill-conditioned $G(x)$, amplifying 
    $G^{-1}$ in~\eqref{eq:fd_christoffel};
    \item Wild $\hat{\Gamma}$ produces large transport norms 
    $\kappa_2 = \|T_{\hat{\Gamma}}\|$, corrupting gradients;
    \item Corrupted gradients prevent $L$ from learning a 
    well-conditioned metric.
\end{enumerate}
Consequently, the condition number proxy 
$\kappa_2(t) = \mathbb{E}[\|\Gamma \cdot \delta\|]$ grows 
monotonically without bound during training.
\end{proposition}

\begin{proof}[Proof sketch]
Each FD evaluation of~\eqref{eq:fd_christoffel} requires $d+1$ 
evaluations of the metric network (one at $x$ and two per 
coordinate direction, with symmetry). With $d = 32$ and batch size 
$B \times T = 1024$ unique contextual positions (no vocabulary 
deduplication possible), this amounts to $33 \times 1024 = 33{,}792$ 
metric network evaluations per step. Each evaluation at a 
previously-unseen contextual position returns an essentially random 
metric tensor until the network has converged --- but convergence 
requires stable gradients, which requires bounded Christoffels, 
closing the loop.

Empirically, on Tiny Shakespeare with $(d, H) = (32, 256)$: 
$\kappa_2$ grows as $54 \to 127 \to 236 \to 322 \to 407 \to 553$ 
over 15 epochs (seed 42, 137), with no sign of stabilisation. 
Training PPL plateaus at $\sim\!15\text{--}16$ versus $V_3$'s 
$1.55$.
\end{proof}

\begin{theorem}[Bounded Transport via Low-Rank Factored 
Connection]\label{thm:lowrank_connection}
Let $R \ll d$ be the \emph{connection rank}. Define the factored 
connection
\begin{equation}\label{eq:lowrank_gamma}
    \tilde{\Gamma}^k_{ij}(x) = \sum_{r=1}^{R} U^k_r(x) \, 
    V_{r,ij}
\end{equation}
where $U(x) \in \mathbb{R}^{d \times R}$ is produced by a 
\emph{connection network} (a learned MLP) and 
$V \in \mathbb{R}^{R \times d \times d}$ is a shared 
(position-independent) parameter. Then:
\begin{enumerate}
    \item \textbf{Bounded transport norms.} The contraction 
    $\tilde{\Gamma} \cdot \delta$ for tangent vector $\delta$ 
    satisfies
    \begin{equation}\label{eq:transport_bound}
        \|\tilde{\Gamma}^k_{ij} \delta^i\|^2 \leq 
        \|U(x)\|_F^2 \cdot \|V\|_F^2 \cdot \|\delta\|^2,
    \end{equation}
    so the transport norm is bounded by the product of the 
    connection-network output norm and the (fixed, learnable) 
    $V$-factor norm. Standard weight initialisation 
    ($V_{r,ij} \sim \mathcal{N}(0, 0.01/\sqrt{R})$) ensures 
    $\|V\|_F = O(d)$ at initialisation, and gradient descent 
    maintains this scale.
    
    \item \textbf{Computational efficiency.} Each transport step 
    requires $O(1)$ connection-network evaluations (versus $O(d)$ 
    metric-network evaluations for FD), reducing the per-step 
    cost from $O(d \cdot C_{\mathrm{metric}})$ to 
    $O(C_{\mathrm{conn}})$ where $C_{\mathrm{metric}}$ and 
    $C_{\mathrm{conn}}$ are the network evaluation costs.
    
    \item \textbf{Rank constraint as regularisation.} The rank-$R$ 
    factorisation restricts $\tilde{\Gamma}$ to an 
    $R$-dimensional subspace of the $d^3$-dimensional space of 
    possible connection coefficients. This implicit regularisation 
    prevents the coarse connection from developing spurious 
    high-frequency structure that the coarse metric cannot resolve.
\end{enumerate}
\end{theorem}

\begin{proof}
(1)~Expanding \eqref{eq:lowrank_gamma}:
\[
    (\tilde{\Gamma} \cdot \delta)^k_j 
    = \sum_i \tilde{\Gamma}^k_{ij} \delta^i 
    = \sum_r U^k_r \underbrace{\left(\sum_i V_{r,ij} 
    \delta^i\right)}_{w_{r,j}}.
\]
By Cauchy--Schwarz in the $r$-index:
$\|(\tilde{\Gamma} \cdot \delta)_j\|^2 
\leq \sum_k \left(\sum_r |U^k_r| |w_{r,j}|\right)^2 
\leq \|U\|_F^2 \|w_j\|^2$,
and $\|w_j\|^2 \leq \|V_{:,:,j}\|_F^2 \|\delta\|^2$ by a second 
application. Summing over $j$ gives \eqref{eq:transport_bound}.

(2)~The FD scheme requires metric evaluations at 
$x \pm h e_i$ for $i = 1, \ldots, d$, giving $2d + 1$ network 
calls. The factored connection~\eqref{eq:lowrank_gamma} requires 
one forward pass of the connection network to produce $U(x)$, 
then matrix operations in $O(dR)$ which are dominated by the 
network evaluation.

(3)~The image of the map $x \mapsto \tilde{\Gamma}(x)$ lies in 
the $R$-dimensional affine subspace 
$\mathrm{span}\{V_1, \ldots, V_R\} \subset 
\mathbb{R}^{d \times d \times d}$. For $R \ll d$, this 
eliminates $(d^3 - Rd^2)$ degrees of freedom, functioning as a 
strong implicit prior toward simple (low-rank) geometry at the 
coarse scale.
\end{proof}

\begin{remark}[Alignment with Asymptotic 
Freedom]\label{rem:lowrank_y6}
The rank constraint on the coarse connection is the architectural 
realisation of Y6 asymptotic freedom. Y6 predicts that coupling 
constants decrease under RG flow --- the effective interaction 
becomes simpler at coarser scales. The factored 
connection~\eqref{eq:lowrank_gamma} makes this structural: the 
fine manifold $\mathcal{M}_1$ uses full finite-difference 
Christoffels (unconstrained connection complexity), while the 
coarse manifold $\mathcal{M}_2$ is restricted to rank-$R$ 
connections. The ratio 
$R / d^2 = 8/1024 \approx 0.008$ measures the degree of 
asymptotic simplification --- the coarse geometry uses less than 
$1\%$ of the connection degrees of freedom available to the fine 
geometry.
\end{remark}

\begin{remark}[Experimental 
Validation]\label{rem:lowrank_experiment}
On Tiny Shakespeare with $(d, H, R) = (32, 256, 8)$:

\medskip
\begin{center}
\begin{tabular}{lcccc}
\toprule
\textbf{Architecture} & \textbf{$\kappa_2$ (ep.\,2)} 
& \textbf{PPL (2\,ep.)} & \textbf{Speed} & \textbf{vs.\,V3} \\
\midrule
V3RG (FD, $R{=}0$) & $54 \to 127$ (diverging) & 27.92 & 
$8.5\times$ slower & $-51\%$ \\
V3RG (Rank-8, $R{=}8$) & $1.2 \to 1.1$ (bounded) & 14.58 & 
$1.9\times$ slower & $+21\%$ \\
V3 (baseline) & --- & 18.46 & $1\times$ & --- \\
\bottomrule
\end{tabular}
\end{center}

\medskip\noindent
The rank-8 connection simultaneously: (i)~bounds $\kappa_2$ near 
unity (versus unbounded growth), (ii)~enables V3RG to outperform 
the V3 baseline from the first evaluation, and 
(iii)~reduces the speed penalty from $8.5\times$ to $1.9\times$ by 
eliminating the $O(d)$ metric-network evaluations. The gate 
parameter remains live at $\sigma(g) \approx 0.49$, confirming 
that the coarse block is actively contributing to the 
representation rather than being gated out.
\end{remark}

% ── Y6c ──
\section{Y6c: Fiber Bundle Transport (Gauge-Coupled Extension)}

The low-rank factored connection (Y6b) stabilises coarse-scale 
transport but introduces a second transport block operating on 
contextual positions---a domain where the learned metric is 
untrained. We show that a \emph{fiber bundle extension} of the 
original transport eliminates this problem entirely: multi-scale 
information is encoded in internal degrees of freedom coupled 
through a gauge connection, within a single transport step over the 
original embedding positions.

\begin{theorem}[Fiber Bundle Transport]\label{thm:fiber_transport}
Let $(\mathcal{M}, g)$ be a Davis manifold with Levi-Civita 
connection $\nabla$ producing Christoffel-contracted transport 
matrices $M(x) \in \mathbb{R}^{d \times d}$ at base points 
$x \in \mathcal{M}$. Let $f \geq 1$ and define the 
\emph{gauge-extended connection}:
\begin{equation}\label{eq:fiber_connection}
    M_{\mathrm{ext}}(x) = 
    \begin{pmatrix} M(x) & W_{bf} \\ 
    -W_{bf}^\top & W_{ff} \end{pmatrix}
    \in \mathbb{R}^{(d+f) \times (d+f)}
\end{equation}
where $W_{bf} \in \mathbb{R}^{d \times f}$ is the 
\emph{gauge coupling} (base $\leftrightarrow$ fiber) and 
$W_{ff} \in \mathbb{R}^{f \times f}$ is the \emph{fiber 
self-connection} (both learned, position-independent). Define the 
extended transport via the Cayley map:
\begin{equation}\label{eq:fiber_cayley}
    R(x) = \bigl(I_{d+f} + \tfrac{1}{2}A(x)\bigr)^{-1}
    \bigl(I_{d+f} - \tfrac{1}{2}A(x)\bigr), \quad
    A(x) = \tfrac{1}{2}\bigl(M_{\mathrm{ext}}(x) - 
    M_{\mathrm{ext}}(x)^\top\bigr).
\end{equation}
Then:
\begin{enumerate}
    \item \textbf{$R(x) \in \SO(d{+}f)$.} The extended transport 
    preserves the norm of the extended state 
    $\tilde{h} = (h_{\mathrm{base}}, h_{\mathrm{fiber}}) 
    \in \mathbb{R}^{d+f}$.
    
    \item \textbf{Multi-scale coupling without metric 
    extrapolation.} The base Christoffels $M(x)$ are evaluated 
    only at original embedding positions $x \in \{e_1, \ldots, 
    e_V\}$, where the metric network is trained. The gauge 
    coupling $W_{bf}$ mixes base and fiber dimensions 
    \emph{algebraically}, without requiring any metric evaluation 
    at new positions. Consequently, the feedback instability of 
    Proposition~\ref{prop:fd_divergence} does not arise.
    
    \item \textbf{The fiber carries holonomy memory.} A state 
    $\tilde{h}_0 = (h_0, f_0)$ transported along a closed loop 
    $\gamma$ on $\mathcal{M}$ returns to 
    $\tilde{h}_0' = R_\gamma \tilde{h}_0$ where 
    $R_\gamma \in \SO(d{+}f)$. The off-diagonal blocks of 
    $R_\gamma$ (generated by $W_{bf}$) transfer holonomy 
    information between base and fiber subspaces: the fiber 
    accumulates geometric information from the base path 
    that persists even when the base component returns to its 
    starting point.
    
    \item \textbf{Dimensional reduction at readout.} Only the 
    first $d$ components of $\tilde{h}$ are used for prediction 
    (logits $= \tilde{h}_{1:d} \cdot E^\top$). The fiber 
    dimensions $\tilde{h}_{d+1:d+f}$ are internal degrees of 
    freedom---analogous to the electromagnetic phase in $U(1)$ 
    gauge theory, which governs dynamics but is not directly 
    observable.
\end{enumerate}
\end{theorem}

\begin{proof}
(1)~$A(x)$ is skew-symmetric by construction, so the Cayley map 
$R = (I + A/2)^{-1}(I - A/2)$ is orthogonal (and 
orientation-preserving since $\det R = +1$ for skew-$A$ in odd and 
even dimensions). 

(2)~The metric network $L(x)$ producing $G(x) = L(x)L(x)^\top + 
\varepsilon I$ is evaluated only at the $V$ vocabulary base points 
(with deduplication), exactly as in the V3 architecture. The gauge 
parameters $W_{bf}, W_{ff}$ are global (position-independent) and 
enter only through the block structure 
of~\eqref{eq:fiber_connection}. No new metric evaluations are 
introduced.

(3)~The holonomy of $\SO(d{+}f)$ transport on a closed loop 
$\gamma$ is $R_\gamma = \prod_{t} R(x_t) \in \SO(d{+}f)$. Write 
$R_\gamma$ in block form:
\[
    R_\gamma = \begin{pmatrix} R_{bb} & R_{bf} \\ R_{fb} & R_{ff}
    \end{pmatrix}.
\]
When $W_{bf} = 0$, the connection is block-diagonal, 
$R_\gamma$ is block-diagonal, and no information transfers 
between base and fiber. For $W_{bf} \neq 0$, the off-diagonal 
blocks $R_{bf}, R_{fb}$ are generically nonzero, enabling 
cross-subspace information flow. The total holonomy group is a 
subgroup of $\SO(d{+}f)$ that is generically \emph{not} 
reducible to $\SO(d) \times \SO(f)$---the gauge coupling generates 
genuinely entangled holonomy.

(4)~The projection $\pi: \mathbb{R}^{d+f} \to \mathbb{R}^d$, 
$\pi(\tilde{h}) = \tilde{h}_{1:d}$, discards the fiber at 
readout. This is the standard construction: in a principal 
$G$-bundle, physical observables are gauge-invariant quantities 
on the base, while the internal fiber degrees of freedom mediate 
dynamics.
\end{proof}

\begin{remark}[Physical Interpretation: Principal Bundle 
Structure]\label{rem:principal_bundle}
The fiber bundle transport has a precise gauge-theoretic 
interpretation. The Davis manifold $(\mathcal{M}, g)$ is the 
\emph{base space}. The fiber $\mathbb{R}^f$ at each point is an 
\emph{internal space} carrying additional geometric degrees of 
freedom. The gauge coupling $W_{bf}$ is the \emph{connection 
form} on the associated vector bundle 
$E = \mathcal{M} \times \mathbb{R}^f$, specifying how the fiber 
directions rotate as the base point moves along a path.

The key insight connecting this to RG: the fiber carries 
\emph{scale-local} information. At each transport step, the base 
subspace encodes the fine-scale (token-level) geometry, while the 
fiber subspace accumulates coarse-scale information through the 
gauge coupling. The coupling $W_{bf}$ plays the role of the 
\emph{anomalous dimension matrix}: it governs how information 
transfers between the fine description (base) and the coarse 
description (fiber) at each step. This is precisely the 
relationship predicted by Y6a (scale-local connections), but 
realised within a \emph{single} transport step rather than by 
stacking separate transport blocks.
\end{remark}

\begin{remark}[Experimental 
Validation]\label{rem:fiber_experiment}
On Tiny Shakespeare with $(d, H, f) = (32, 256, 8)$, five 
physics-inspired architectures were tested at iso-parameters 
($\sim\!154$K each), two seeds $\{42, 137\}$, five epochs:

\medskip
\begin{center}
\begin{tabular}{lcccl}
\toprule
\textbf{Architecture} & \textbf{PPL ($\pm$ std)} 
& \textbf{vs.\ V3} & \textbf{Time} & \textbf{Physics} \\
\midrule
V3-Holo (AdS/CFT) & 37.45 & $+217\%$ & $2.9\times$ 
& two transports, shared $\mathcal{M}$ \\
V3-BO (Born--Opp.) & 34.85 & $+195\%$ & $7.5\times$ 
& $\kappa$-conditioned $T_2$ \\
\midrule
V3 (baseline) & $11.82 \pm 0.00$ & --- & $1\times$ & --- \\
V3-Y2 (multi-head) & $11.09 \pm 0.20$ & $-6.2\%$ & $3.0\times$ 
& block-diagonal sub-transports \\
\textbf{V3-Fiber (gauge)} & $\mathbf{11.06 \pm 0.02}$ 
& $\mathbf{-6.5\%}$ & $\mathbf{1.4\times}$ 
& \textbf{$\SO(d{+}f)$ with $W_{bf}$} \\
\bottomrule
\end{tabular}
\end{center}

\medskip\noindent
Three findings:
\begin{enumerate}[nosep]
    \item \emph{Two-transport architectures fail catastrophically.} 
    Holo and BO both evaluate the learned metric at post-attention 
    positions where it is untrained, producing wild Christoffels 
    ($\kappa_2 = 8$--$10$). This confirms 
    Proposition~\ref{prop:fd_divergence}: the instability is not 
    specific to finite-difference estimation but to any metric 
    evaluation at out-of-distribution positions.
    
    \item \emph{Fiber transport wins on all axes.} Best PPL 
    ($11.06$), lowest variance ($\pm 0.02$, compared to Y2's 
    $\pm 0.20$), and minimal overhead ($1.4\times$, compared to 
    Y2's $3.0\times$). The gauge coupling converges to 
    $\|W_{bf}\| \approx 0.13$ with consistent coupling energy 
    $\approx 0.28$---small but nonzero, confirming that the 
    fiber is active.
    
    \item \emph{The universal physics principle holds.} Every 
    architecture where coarse transport emerges from or is coupled 
    to the fine geometry (Fiber, Y2) beats V3. Every architecture 
    where the metric is evaluated at unfamiliar positions (Holo, 
    BO) or where the coarse geometry is independent (the earlier 
    V3RG) fails. This is the computational confirmation of the 
    gauge-theoretic prediction of Y6a.
\end{enumerate}
\end{remark}

% ── Y7 ──
\section{Y7: Partition Function Decimation}

\begin{conjecture}[Decimation 
Formula]\label{conj:decimation}
The coarsened partition function $Z_0(\beta, \ell)$ is related to 
the original by a \emph{decimation formula}:
\begin{equation}\label{eq:decimation}
    Z_0(\beta, \ell) = \sum_{\sigma_\ell \in 
    \Cech^0(\nerve_\ell, \sheaf_\ell)} 
    \exp(-\beta \, H_{\mathrm{eff}}(\sigma_\ell))
\end{equation}
where the effective Hamiltonian is defined by integrating out 
sub-scale degrees of freedom:
\begin{equation}\label{eq:heff}
    \exp(-\beta \, H_{\mathrm{eff}}(\sigma_\ell)) 
    = \sum_{\substack{\sigma \in \Cech^0(\nerve, \sheaf) \\ 
    \pi_\ell(\sigma) = \sigma_\ell}} 
    \exp(-\beta \, \langle \sigma, \Delta_0 \, \sigma \rangle)
\end{equation}
The sum runs over all fine-scale 0-cochains $\sigma$ consistent 
with the coarse-scale configuration $\sigma_\ell$.

\textbf{Normalization.} The decimation preserves the total partition 
function: $Z_0(\beta, 0) = Z_0(\beta, \ell)$ since every fine-scale 
configuration $\sigma$ belongs to exactly one coarse-scale class 
$\sigma_\ell = \pi_\ell(\sigma)$. No Jacobian or measure factor is 
needed because the summation variable changes from fine-scale 
configurations to (coarse-scale configuration, sub-scale 
configuration within block) pairs, and the latter decomposition is 
a partition of the original sum. When passing to continuous 
integration (for large stalks), a Jacobian from the change of 
variables would appear; in the discrete setting used here, 
combinatorial counting suffices.

The effective Hamiltonian $H_{\mathrm{eff}}$ generates new 
``interactions'' (non-local terms) at each RG step --- the standard 
phenomenon of RG-generated couplings. Notably, $H_{\mathrm{eff}}$ 
is no longer purely quadratic even if the original Hamiltonian 
$\langle \sigma, \Delta_0 \, \sigma \rangle$ was, reflecting the 
nonlinearity of constraint satisfaction.
\end{conjecture}

\begin{remark}[Truncation and Asymptotic 
Freedom]\label{rem:truncation}
The \v{C}ech Laplacian quadratic form is not closed under RG: after 
one decimation step, $H_{\mathrm{eff}}$ contains quartic and 
higher-order terms. The \v{C}ech framework provides a 
\emph{leading-order (Gaussian) approximation} to the exact effective 
theory. By Y6 (asymptotic freedom), the generated higher-order 
couplings are suppressed at large scales --- the Gaussian 
approximation improves under coarse-graining. Rigorous control of 
the remainder requires bounding the generated higher-order 
couplings, which we leave as an open problem.
\end{remark}

\textbf{Derived from:} RG8 (monotonicity) + Branch~VII, SH7 
(partition function).

% ── Y8 ──
\section{Y8: Phase Diagram Flow Lines}

\begin{conjecture}[RG Flow on the Phase 
Diagram]\label{conj:flow_lines}
The phase diagram of the parent framework (Conjecture~50.1) in the 
$(\rho, \Khat)$ plane acquires RG flow lines. Each initial 
condition $(\rho_0, \Khat_0)$ flows under $T_\ell$ to a trajectory 
$(\rho(\ell), \Khat(\ell))$:

\begin{enumerate}[label=(\alph*)]
    \item Above the critical line: flow to the determined fixed 
    point ($\rho \to \infty$, $\Khat \to 0$).
    
    \item Below the critical line: flow to the underdetermined 
    fixed point ($\rho \to 0$, $\Khat \to 0$).
    
    \item On the critical line: flow to the critical fixed point 
    $(\rho^*, \Khat^*)$.
\end{enumerate}

The critical line $\rho_{\mathrm{crit}}(\Khat)$ is the 
\emph{separatrix} of the flow (the invariant manifold separating 
the two basins of attraction), and its equation is determined by 
$\beta(\Gamma) = 0$, i.e., the condition $\Gamma = 1$.
\end{conjecture}

\textbf{Derived from:} RG4--RG7 + Parent Conjecture~50.1.

% ── Y9 ──
\section{Y9: Finite-Size Scaling}

\begin{conjecture}[Finite-Size 
Effects]\label{conj:finite_size}
For a completion problem on a finite nerve with $N$ vertices 
(e.g., Sudoku with $N = 27$ constraints), the RG flow terminates 
at scale $\ell_{\max} \sim N^{1/d_H}$ (the system size). 
Finite-size effects replace the true phase transition with a 
smooth crossover.

The \emph{finite-size scaling ansatz}:
\begin{equation}\label{eq:fss}
    \Gamma_{\mathrm{eff}}(N) - 1 
    \sim N^{-1/(\nu d_H)} \cdot 
    G\!\left(N^{1/(\nu d_H)} \cdot (\Gamma_0 - 1)\right)
\end{equation}
where $G$ is a universal scaling function. This predicts:
\begin{align}
    \text{Critical region width:}\quad 
    \Delta\mu &\sim N^{-1/(\nu d_H)} 
    \label{eq:fss_width}\\
    \text{Maximum susceptibility:}\quad 
    \chi_{\max} &\sim N^{\gamma/(\nu d_H)} 
    \label{eq:fss_chi}\\
    \text{Order parameter at criticality:}\quad 
    \langle m \rangle &\sim N^{-\beta_{\mathrm{mag}}/(\nu d_H)} 
    \label{eq:fss_m}
\end{align}

\textbf{Testable prediction for Sudoku:} A $4 \times 4$ Shidoku 
($N = 12$ constraints) has a much broader crossover than a 
$9 \times 9$ Sudoku ($N = 27$) or $16 \times 16$ super-Sudoku 
($N = 48$). The sharpness of the determined/underdetermined 
transition scales as $N^{1/(\nu d_H)}$.
\end{conjecture}

\textbf{Derived from:} RG5 (fixed-point structure) + RG6 (critical 
exponents) + Branch~VII, Appendix~A (Shidoku computation).

% ── Y10 ──
\section{Y10: Semantic Depth as RG Scale}

\begin{conjecture}[Semantic 
Depth]\label{conj:semantic_depth}
The ``depth'' of a semantic completion --- the number of reasoning 
steps required --- equals the number of RG steps needed to reach 
the determined fixed point:
\begin{equation}\label{eq:semantic_depth}
    \mathrm{depth} = \min\{\ell : 
    \Gamma(\ell) > \Gamma_{\mathrm{threshold}}\}
\end{equation}
This provides a precise, mathematically grounded definition of 
\emph{semantic depth} connecting three quantities:

\begin{enumerate}[label=(\roman*)]
    \item The mathematical RG (number of coarse-graining steps on 
    the nerve).
    \item The computational architecture (number of transformer 
    layers actively engaged).
    \item The cognitive difficulty (number of reasoning steps a 
    human needs).
\end{enumerate}

Shallow completions (syntax, pattern matching) have 
$\mathrm{depth} \sim 1$: one coarse-graining step suffices. 
Deep completions (multi-step reasoning, long-range inference) have 
$\mathrm{depth} \sim L$: the full transformer depth is needed. The 
maximum semantic depth of a model is bounded by its architectural 
depth $L$.
\end{conjecture}

\textbf{Derived from:} Y1 (layers as RG steps) + Y3 (completion 
capacity).

% ═══════════════════════════════════════════════════════
\newpage
\part{Open Questions}
% ═══════════════════════════════════════════════════════

\section{Directions for Future Work}

\begin{enumerate}
    \item \textbf{Explicit beta function computation:} The 
    qualitative form $\beta(\Gamma) = (1/\nu) \cdot \Gamma \cdot 
    (\Gamma - 1) \cdot h(\Gamma)$ is constrained by fixed-point 
    structure, flow direction, and Branch~VI consistency. Computing 
    $h(\Gamma)$ for specific completion problems (Shidoku, random 
    3-SAT) via numerical RG on the nerve is the most direct test.
    
    \item \textbf{Upper critical dimension:} Is 
    $d_{\mathrm{uc}} = 4$ for the Davis completion transition, as 
    for the Ising model? Or does the non-standard interaction 
    structure of constraint graphs change the critical dimension? 
    Numerical study of completion transitions on lattices of varying 
    dimension could answer this.
    
    \item \textbf{Universality class identification:} Do Sudoku and 
    Latin squares share the same critical exponents? Do random 
    $k$-SAT and Erd\H{o}s--R\'enyi graph coloring? Testing 
    universality requires computing $\nu, \alpha, \eta$ for 
    specific problem families and comparing.
    
    \item \textbf{Transformer mechanistic interpretation:} The 
    layer-scale correspondence (Y1) predicts that per-layer 
    activation norms peak at critical layers and that attention 
    patterns become sparse after convergence. These predictions are 
    testable on existing transformer architectures via 
    mechanistic interpretability methods.
    
    \item \textbf{Rigorous C-theorem:} 
    Theorem~\ref{thm:c_theorem} proves monotonicity of the trace. 
    The universality of $C_{\mathrm{crit}}$ (the value of 
    completion entropy at the critical fixed point) remains 
    conjectural. Is $C_{\mathrm{crit}}$ computable from the 
    universality class data ($d_H$, symmetry group, range)?
    
    \item \textbf{Anomalous dimension from BCH:} The perturbative 
    formula $\eta = c \cdot \|\alpha_c\|^2$ (RG2a) requires 
    computing the BCH correction to the two-point function at 
    criticality. Is the perturbative expansion convergent, or does 
    a Wilson--Fisher-type fixed point require non-perturbative 
    methods?
    
    \item \textbf{Asymptotic freedom beyond leading order:} Y6 
    establishes that the linearized theory improves at large scales. 
    What is the rate of improvement? Can we bound the generated 
    higher-order couplings (Y7, Remark~\ref{rem:truncation}) to 
    establish rigorous control of the Gaussian approximation?
    
    \item \textbf{Multi-scale persistence and computational 
    complexity:} The obstruction persistence diagram (Y5) 
    classifies completion difficulty by RG scale. Is there a 
    relationship between the persistence diagram and the 
    computational complexity (NP-hardness, approximation ratios) 
    of the associated CSP?
\end{enumerate}

% ═══════════════════════════════════════════════════════
\newpage
\appendix
\part*{Appendix}
\addcontentsline{toc}{part}{Appendix}
% ═══════════════════════════════════════════════════════

\section{Finite-Model Anchor: RG Flow on 
Shidoku}\label{app:shidoku_rg}

To ground the abstract machinery in a computable example, we carry 
out the full RG flow for the $4 \times 4$ Shidoku grid from the 
Branch~VII appendix, verifying that the coarse-graining 
construction, Cauchy interlace inequality, and Euler characteristic 
invariance hold explicitly.

\subsection{Setup: The Shidoku Nerve}

The Shidoku nerve $\nerve(\cover)$ has:
\begin{itemize}
    \item 12 vertices: $r_1, r_2, r_3, r_4$ (rows), 
    $c_1, c_2, c_3, c_4$ (columns), 
    $b_1, b_2, b_3, b_4$ (boxes).
    \item 32 edges: 16 row-column (each sharing 1 cell), 
    8 row-box (each sharing 2 cells), 8 column-box (each sharing 
    2 cells).
    \item 16 triangles: one per cell (triple overlap of row, 
    column, and box).
\end{itemize}

The adjacency structure of the row-box and column-box overlaps is:
\begin{align*}
    r_1, r_2 &\text{ overlap } b_1, b_2 
    && c_1, c_2 \text{ overlap } b_1, b_3 \\
    r_3, r_4 &\text{ overlap } b_3, b_4 
    && c_3, c_4 \text{ overlap } b_2, b_4
\end{align*}
No two rows share cells. No two columns share cells. No two boxes 
share cells.

\subsection{Scale $\ell = 0$: Original Problem}

At scale $\ell = 0$, $\nerve_0 = \nerve$ with 12 vertices and 32 
edges. The graph Laplacian $L_0 = D - A$ (where $D$ is the degree 
matrix and $A$ the adjacency matrix) on scalar 0-cochains has 
spectrum:
\begin{equation}\label{eq:shidoku_spectrum_0}
    \mathrm{spec}(L_0) = \{0,\; 2.764,\; 2.764,\; 4,\; 6,\; 6,\; 
    6,\; 6,\; 6,\; 7.236,\; 7.236,\; 10\}
\end{equation}
The spectral gap is $\lambda_1 = 4(2 - \sqrt{2}) \approx 2.764$.

\subsection{Scale $\ell = 1$: Band-Level Coarse-Graining}

The natural coarse-graining groups constraints into bands:
\begin{align*}
    \block_1 &= \{r_1, r_2, b_1, b_2\} 
    && \text{(top band: 2 rows + 2 boxes)} \\
    \block_2 &= \{r_3, r_4, b_3, b_4\} 
    && \text{(bottom band)} \\
    \block_3 &= \{c_1, c_2\} && \text{(left columns)} \\
    \block_4 &= \{c_3, c_4\} && \text{(right columns)}
\end{align*}
This partitions all 12 nerve vertices. The blocks satisfy 
$\diam(\block_a) \leq 1$ in $d_\nerve$ (every vertex in a block 
is adjacent to every other vertex in that block, since rows and 
boxes in the same band share cells).

The coarsened nerve $\nerve_1$ has 4 vertices. An edge 
$[\block_a, \block_b]$ exists whenever some $v \in \block_a$ is 
adjacent to some $w \in \block_b$ in $\nerve$. Since every row 
overlaps every column but no column overlaps another column and 
no row-band overlaps another row-band:
\begin{equation}\label{eq:coarsened_edges}
    \text{Edges of } \nerve_1 = \{(\block_1, \block_3),\; 
    (\block_1, \block_4),\; (\block_2, \block_3),\; 
    (\block_2, \block_4)\}
\end{equation}
Thus $\nerve_1 \cong K_{2,2}$ (the complete bipartite graph).

\subsection{Laplacian at Scale $\ell = 1$}

With the inherited inner product (block-compatible 0-cochains 
embedded in $\mathbb{R}^{12}$, with each block's indicator 
normalized to unit length), the \emph{compression} 
$\Delta_0^{(1)} = P_1^\dagger L_0 P_1$ 
(Definition~\ref{def:coarsened_cech}) has spectrum:
\begin{equation}\label{eq:shidoku_spectrum_1}
    \mathrm{spec}(\Delta_0^{(1)}) = \{0,\; 3,\; 6,\; 9\}
\end{equation}
The spectral gap increases: $\lambda_1 = 2.764 \to 3.000$. (By 
contrast, the bare graph Laplacian of $K_{2,2}$ with unit edge 
weights has spectrum $\{0, 2, 2, 4\}$ --- see 
Remark~\ref{rem:compression_vs_coarse}.)

\subsection{Verification}

\textbf{(1) Cauchy interlace.} Sorting both spectra and comparing 
the first 4 eigenvalues:
\begin{center}
\renewcommand{\arraystretch}{1.2}
\begin{tabular}{cccl}
$j$ & $\lambda_j$ (original) & $\mu_j$ (coarsened) 
& Interlace \\
\hline
0 & $0$ & $0$ & $\mu_0 \geq \lambda_0$ \checkmark \\
1 & $2.764$ & $3$ & $\mu_1 \geq \lambda_1$ \checkmark \\
2 & $2.764$ & $6$ & $\mu_2 \geq \lambda_2$ \checkmark \\
3 & $4$ & $9$ & $\mu_3 \geq \lambda_3$ \checkmark \\
\end{tabular}
\end{center}

\textbf{(2) Trace inequality.} At $\beta = 1$:
\begin{align}
    \tr e^{-\beta L_0} &= \sum_{j=0}^{11} 
    e^{-\lambda_j} \approx 1.158 
    \label{eq:trace_full}\\
    \tr e^{-\beta \Delta_0^{(1)}} &= 1 + e^{-3} + e^{-6} 
    + e^{-9} \approx 1.052 
    \label{eq:trace_coarse}
\end{align}
confirming $C(1) = \ln(1.052) \leq C(0) = \ln(1.158)$: completion 
entropy decreases under coarse-graining.

\textbf{(3) Euler characteristic invariance.} For scalar 
0-cochains on the nerve (the simplified setting of this appendix), 
the Euler characteristic $\chi = \dim \ker L - \dim \Hcech^1 
+ \cdots$ The zero eigenvalue is preserved 
($\lambda_0 = \mu_0 = 0$), and the alternating sum of Betti 
numbers is invariant by the Leray spectral sequence argument 
(Theorem~\ref{thm:euler_invariant}).

\textbf{(4) Flow direction.} The spectral gap \emph{increases} 
under coarse-graining ($2.764 \to 3.000$), corresponding to 
$m(\ell)^2$ increasing and the effective correlation length 
$\xi(\ell) = 1/m(\ell)$ decreasing. For a well-posed Shidoku 
(unique solution, $\Gamma > 1$), this confirms the RG flow toward 
the determined fixed point: constraints become easier to satisfy at 
coarser scales.

\subsection{Three Shidoku Cases}

For the three cases from Branch~VII, Appendix~A:

\begin{enumerate}
    \item \textbf{Well-posed (6 clues, unique solution):} Clues 
    reduce stalk dimensions. At scale $\ell = 0$: 
    $\dim \Hcech^0 = 1$, $\Hcech^1 = 0$. Under coarse-graining, 
    the unique solution projects to a unique coarsened section. 
    $\Gamma(\ell)$ increases: flow to determined fixed point.
    
    \item \textbf{Under-clued (2 clues, 48 solutions):} 
    $\dim \Hcech^0 = 48$, $\Hcech^1 = 0$. Coarse-graining reduces 
    the number of distinguishable solutions (some differ only at 
    sub-scale level). $\Gamma(\ell)$ decreases: flow toward 
    underdetermined fixed point.
    
    \item \textbf{Contradictory (locally valid, globally 
    impossible):} $\Hcech^0 = 0$, $\dim \Hcech^1 \geq 1$. The 
    obstruction class $[\alpha] \in \Hcech^1$ is either relevant 
    (persists at scale $\ell = 1$, detectable from band-level 
    incompatibility) or irrelevant (vanishes at scale $\ell = 1$, 
    a sub-band obstruction). Under Leray-acyclic coarse-graining 
    (Theorem~\ref{thm:euler_invariant}), the Euler characteristic 
    $\chi(\sheaf) = 0 - \dim \Hcech^1 < 0$ is preserved at all 
    scales.
\end{enumerate}

% ═══════════════════════════════════════════════════════
\newpage
\section{Result Summary}\label{app:results}

\begin{center}
\renewcommand{\arraystretch}{1.3}
\begin{tabular}{llcl}
\textbf{Label} & \textbf{Title} & \textbf{Type} 
& \textbf{Dependencies} \\
\hline
\multicolumn{4}{c}{\emph{Foundational (RG1--RG8)}} \\
\hline
RG1 & Block-Spin Coarse-Graining & Def + Thm 
& VII.SH1, VII.SH2 \\
RG2 & Coarse-Graining Semigroup & Def 
& RG1 \\
RG3 & \v{C}ech Laplacian at Scale $\ell$ & Def + Conj 
& VII.SH7, VII.X3 \\
RG4 & Beta Function & Conj 
& RG2, RG3, VI.S7 \\
RG5 & Fixed Points and Stability & Conj 
& RG4, VI.S7 \\
RG6 & Critical Exponents from RG & Conj 
& RG5, VI.S4 \\
RG7 & Phase Diagram as RG Flow & Conj 
& RG4--5, Parent 50.1 \\
RG8 & Monotonicity (C-Theorem) & Thm 
& RG2, VII.SH7 \\
\hline
\multicolumn{4}{c}{\emph{First-Order Derivations 
(RG1a--RG7a)}} \\
\hline
RG1a & Derived Pushforward / Spectral Seq.\ & Thm + Conj 
& RG1, VII.SH5a \\
RG2a & Anomalous Dimensions & Conj 
& RG6, VII.SH5 \\
RG3a & Crossover Scaling Functions & Conj 
& RG5, RG6, VI \\
RG4a & Corrections to Scaling & Conj 
& RG5, RG6 \\
RG5a & Universality Classes & Conj 
& RG6 \\
RG6a & RG-Improved Partition Function & Conj 
& RG2, RG4, VII.SH7 \\
RG7a & Euler Characteristic Invariance & Thm 
& RG1a, VII.SH7 \\
\hline
\multicolumn{4}{c}{\emph{Cross-Synthesis (Y1--Y10, Y6a--b)}} \\
\hline
Y1 & Layers as RG Steps & Conj 
& RG4, RG5 \\
Y2 & Multi-Head = Multi-Scale & Conj 
& Y1, VII.SH5 \\
Y3 & Completion Capacity & Conj 
& Y1, RG5, RG6 \\
Y4 & Residual Stream = Trajectory & Conj 
& Y1 \\
Y5 & Persistence = Relevance & Conj 
& RG1a, VII.SH4b \\
Y6 & Asymptotic Freedom & Conj 
& VII.SH5, RG2 \\
Y6a & Scale-Local Connections & Conj 
& Y1, RG3, Y6 \\
Y6b & Bounded Coarse Transport & Prop+Thm 
& Y6a, Y6 \\
Y6c & Fiber Bundle Transport & Thm 
& Y6a, Y6b \\
Y7 & Partition Function Decimation & Conj 
& RG8, VII.SH7 \\
Y8 & Phase Diagram Flow Lines & Conj 
& RG4--7 \\
Y9 & Finite-Size Scaling & Conj 
& RG5, RG6 \\
Y10 & Semantic Depth = RG Scale & Conj 
& Y1, Y3 \\
\end{tabular}
\end{center}

\medskip\noindent
\textbf{Totals:} 29 results (6~theorems, 1~proposition, 22~conjectures). 
Framework total after Branch~VIII: $89 + 44 + 30 + 29 = 192$ 
results.

\end{document}
